\documentclass[11pt,aspectratio=169]{beamer}

\usepackage[utf8]{inputenc}
\usepackage[T1]{fontenc}
\usepackage[french]{babel}
\shorthandoff{:}
\usepackage{csquotes}
\usepackage{libertinus}
\usepackage{microtype}
\usepackage{amsmath,amssymb,mathtools}
\usepackage{booktabs,tabularx,array,colortbl}
\usepackage{graphicx}
\makeatletter
% Les sources longues portent des descriptions alternatives PDF. Beamer ne
% prend pas encore en charge DocumentMetadata : la clé reste donc sans effet
% dans le diaporama, sans qu'il faille dupliquer les appels d'image.
\define@key{Gin}{alt}{}
\makeatother
\usepackage[export]{adjustbox}
\usepackage{tikz}
\usetikzlibrary{arrows.meta,positioning,calc,decorations.pathreplacing,fit,backgrounds,patterns}
\tikzset{alt/.code={}}
\usepackage{pgfplots}
\usepgfplotslibrary{statistics}
\pgfplotsset{compat=1.18}
\usepackage{xcolor}
\usepackage{listings}

\definecolor{ink}{HTML}{18212B}
\definecolor{accent}{HTML}{155E75}
\definecolor{accentlight}{HTML}{E7F3F6}
\definecolor{warning}{HTML}{9A3412}
\definecolor{warninglight}{HTML}{FFF1E8}
\definecolor{gridgray}{HTML}{CBD5E1}
\definecolor{rulecolor}{HTML}{4D7C0F}
\definecolor{rulelight}{HTML}{F1F8E8}
\definecolor{courseproblem}{HTML}{687386}
\definecolor{coursedata}{HTML}{47777B}
\definecolor{coursescore}{HTML}{356B8B}
\definecolor{courseprobability}{HTML}{5B6093}
\definecolor{courseloss}{HTML}{8A5D68}
\definecolor{coursegradient}{HTML}{98633F}
\definecolor{courseupdate}{HTML}{8B7040}
\definecolor{coursedecision}{HTML}{687548}
\definecolor{courseevaluation}{HTML}{527460}

% Les deux matières reçoivent une teinte propre. Les maisons occupant le bleu
% d'ardoise (coursescore, 204°) et la terre de Sienne (warning, 19°), Potions
% prend le violet (271°) et Vol le vert (123°) : les quatre teintes sont
% séparées d'au moins 67° et ne se confondent à aucune taille de corps.
\definecolor{coursepotions}{HTML}{6D3F9E}
\definecolor{coursevol}{HTML}{2E7D32}

% Grille 16:9 : toute la surface sert au cours, sans chrome de navigation.
\renewcommand{\familydefault}{\sfdefault}
\usefonttheme{professionalfonts}
\newlength{\coursesidemargin}
\setlength{\coursesidemargin}{7mm}
\setbeamersize{text margin left=\coursesidemargin,text margin right=\coursesidemargin}
\setbeamertemplate{navigation symbols}{}
\setbeamertemplate{headline}{}
\setbeamertemplate{footline}{}
\setbeamertemplate{caption}[numbered]
\setbeamertemplate{frametitle continuation}{%
  \ifnum\insertcontinuationcount>1
    \space{\small\mdseries\color{white!62}(suite)}%
  \fi}
\setbeamercolor{normal text}{fg=ink,bg=white}
\setbeamercolor{structure}{fg=accent}
\setbeamercolor{frametitle}{fg=white,bg=ink}
\setbeamercolor{caption name}{fg=accent!72!ink}
\setbeamercolor{caption}{fg=ink!62}
\setbeamerfont{normal text}{size=\normalsize}
\setbeamerfont{frametitle}{size=\footnotesize,series=\bfseries}
\setbeamerfont{framecrumb}{size=\tiny,series=\bfseries}
\setbeamerfont{caption}{size=\scriptsize}
\setbeamerfont{section title}{size=\LARGE,series=\bfseries}
\setbeamerfont{part title}{size=\LARGE,series=\bfseries}
\setbeamertemplate{frametitle}{%
  \nointerlineskip
  \begin{tikzpicture}[remember picture,overlay]
    \fill[ink] (current page.north west)
      rectangle ([yshift=-5mm]current page.north east);
    \node[anchor=center,inner sep=0pt,outer sep=0pt]
      at ([yshift=-2.5mm]current page.north) {%
        \begin{adjustbox}{max width=.84\paperwidth}%
          {\usebeamerfont{frametitle}\color{white}%
            \MakeUppercase{\insertframetitle}}%
        \end{adjustbox}};
  \end{tikzpicture}%
  \rule{0pt}{5mm}\par
  \nointerlineskip
  \coursenavigation}
\setbeamertemplate{itemize item}{\color{accent!62}\small\textbullet}
\setbeamertemplate{itemize subitem}{\color{accent!48}\scriptsize\textbullet}
\setbeamertemplate{enumerate item}{%
  \colorbox{accent!12}{\color{accent!82!ink}\scriptsize\bfseries\insertenumlabel}}
\setlength{\parskip}{.42em}
\setlength{\emergencystretch}{2em}
\linespread{1.05}

\title{Régression logistique binaire}
\subtitle{Du score linéaire à la décision probabiliste}
\author{Sylvain Maitre}
\date{}
\setbeamertemplate{title page}{%
  \centering
  \vfill
  {\fontsize{29}{33}\selectfont\bfseries\color{ink}\inserttitle\par}
  \vspace{3mm}
  {\fontsize{16}{20}\selectfont\color{accent}\insertsubtitle\par}
  \vspace{7mm}{\color{accent!42}\rule{31mm}{1pt}}\par
  \vfill
  {\small\color{ink!72}\insertauthor\par}
  \vspace{4mm}}

% Les images respectent simultanément la largeur disponible et la hauteur utile.
\let\courseincludegraphics\includegraphics
\RenewDocumentCommand{\includegraphics}{O{}m}{%
  \courseincludegraphics[max height=.67\textheight,keepaspectratio,#1]{#2}}
\NewDocumentEnvironment{slidevisual}{}{%
  \begin{adjustbox}{max totalsize={.97\textwidth}{.75\textheight},center}%
  \begin{minipage}{\textwidth}}
  {\end{minipage}\end{adjustbox}}

\pgfplotsset{courseplot/.style={width=.76\textwidth,height=5.25cm,
  grid=major,grid style={gridgray!40}}}
\pgfplotsset{cloud/.style={width=.72\textwidth,height=5.4cm,
  axis lines=box,axis line style={ink!35},
  grid=major,grid style={gridgray!35},
  tick label style={font=\scriptsize},label style={font=\scriptsize},
  xlabel style={yshift=-2pt},ylabel style={yshift=4pt},
  legend style={at={(0.5,1.08)},anchor=south,legend columns=-1,draw=none,
    font=\scriptsize,/tikz/every even column/.append style={column sep=8pt}}}}
\pgfplotsset{panel/.style={width=.3\textwidth,height=4.15cm,
  axis lines=box,axis line style={ink!35},
  grid=major,grid style={gridgray!35},
  tick label style={font=\tiny},label style={font=\tiny},
  title style={font=\scriptsize,yshift=-2pt},
  xmin=-2.8,xmax=2.8,ymin=-2.8,ymax=2.8,
  xtick={-2,0,2},ytick={-2,0,2}}}

\lstset{
  basicstyle=\ttfamily\scriptsize,
  keywordstyle=\color{accent}\bfseries,
  commentstyle=\color{rulecolor},
  stringstyle=\color{warning},
  showstringspaces=false,
  breaklines=true,
  breakatwhitespace=true,
  frame=none,
  numbers=none,
  columns=fullflexible,
  language=Python,
  literate={é}{{\'e}}1 {è}{{\`e}}1 {ê}{{\^e}}1
    {à}{{\`a}}1 {ç}{{\c c}}1}
\lstdefinestyle{code}{
  basicstyle=\ttfamily\scriptsize,
  xleftmargin=3mm,
  aboveskip=.5em,
  belowskip=.4em}

% Numérotation stable entre le cours long et le diaporama.
\newcounter{coursechapter}
\newcounter{coursestep}
\newcommand{\coursechapterkind}{Chapitre}
\newif\ifcoursereperemode
\coursereperemodefalse
\newcommand{\coursechapterindicator}{}
\newcommand{\coursepartshort}{Vue d'ensemble}
\newcommand{\coursepartposition}{Vue d'ensemble}
\newcommand{\coursepartindex}{0}
\newcommand{\coursepartcolor}{courseproblem}
\newif\ifcoursenavigationvisible
\coursenavigationvisiblefalse
\newcount\courseslideinstage
\newcount\coursestagetotal
\newcount\coursenextstage
\newdimen\courseprogresswidth
% La progression est calculée exclusivement par LaTeX à partir des numéros de
% frames enregistrés dans le fichier auxiliaire ; aucun total n'est généré.
\newcommand{\coursefinishstage}{}
\newcommand{\courseprogressdimensions}{%
  \courseslideinstage=1
  \coursestagetotal=1
  \ifnum\getrefnumber{course-stage-start-\coursepartindex}>0
    \courseslideinstage=\numexpr\insertframenumber-
      \getrefnumber{course-stage-start-\coursepartindex}+1\relax
    \coursenextstage=\numexpr\coursepartindex+1\relax
    \ifnum\coursenextstage>8
      \coursestagetotal=\numexpr\inserttotalframenumber-
        \getrefnumber{course-stage-start-\coursepartindex}+1\relax
    \else
      \ifnum\getrefnumber{course-stage-start-\the\coursenextstage}>0
        \coursestagetotal=\numexpr
          \getrefnumber{course-stage-start-\the\coursenextstage}-
          \getrefnumber{course-stage-start-\coursepartindex}\relax
      \fi
    \fi
  \fi
  \ifnum\courseslideinstage<1
    \courseslideinstage=1
  \fi
  \ifnum\courseslideinstage>\coursestagetotal
    \courseslideinstage=\coursestagetotal
  \fi
  \courseprogresswidth=\dimexpr
    .125\paperwidth*\courseslideinstage/\coursestagetotal\relax}
\newcommand{\coursestage}[3]{%
  \begingroup
  \ifnum#1=\coursepartindex
    \courseprogressdimensions
    \def\stageforeground{ink}%
    \def\stagebase{#3!14}%
    \def\stageweight{\bfseries}%
    \def\stagefillwidth{\courseprogresswidth}%
    \def\stagefillcolor{#3!48}%
  \else
    \ifnum#1<\coursepartindex
      \def\stageforeground{white}%
      \def\stagebase{#3!78!ink}%
      \def\stagefillwidth{0pt}%
      \def\stagefillcolor{#3!78!ink}%
    \else
      \def\stageforeground{#3!78!ink}%
      \def\stagebase{#3!10}%
      \def\stagefillwidth{0pt}%
      \def\stagefillcolor{#3!10}%
    \fi
    \def\stageweight{\mdseries}%
  \fi
  \pgfmathsetlengthmacro{\stageleft}{(#1-1)*.125*\paperwidth}%
  \pgfmathsetlengthmacro{\stagecenter}{\stageleft+.0625*\paperwidth}%
  \fill[\stagebase]
    (\stageleft,0) rectangle ++(.125\paperwidth,-3.8mm);
  \fill[\stagefillcolor]
    (\stageleft,0) rectangle ++(\stagefillwidth,-3.8mm);
  \draw[white,line width=.8pt]
    (\stageleft,0) rectangle ++(.125\paperwidth,-3.8mm);
  \node[anchor=center,inner sep=0pt,outer sep=0pt]
    at (\stagecenter,-1.9mm) {%
      \fontsize{5.55}{6}\selectfont\stageweight
      \color{\stageforeground}\MakeUppercase{#2}};
  \endgroup}
\newcommand{\coursenavigation}{%
  \ifcoursenavigationvisible
    \begingroup
    \leftskip=0pt\rightskip=0pt\parfillskip=0pt
    \begin{tikzpicture}[remember picture,overlay]
      \begin{scope}[shift={([yshift=-5mm]current page.north west)}]
        \coursestage{1}{Données}{coursedata}%
        \coursestage{2}{Score linéaire}{coursescore}%
        \coursestage{3}{Sigmoïde}{courseprobability}%
        \coursestage{4}{Perte logistique}{courseloss}%
        \coursestage{5}{Gradient}{coursegradient}%
        \coursestage{6}{Mise à jour}{courseupdate}%
        \coursestage{7}{Décision}{coursedecision}%
        \coursestage{8}{Évaluation}{courseevaluation}%
      \end{scope}
    \end{tikzpicture}%
    \rule{0pt}{3.8mm}%
    \endgroup
  \fi}
\renewcommand{\thecoursechapter}{\arabic{coursechapter}}
\newcommand{\courseintro}{%
  \coursefinishstage
  \coursenavigationvisiblefalse
  \gdef\coursechapterindicator{}%
  \gdef\coursepartshort{Introduction}%
  \gdef\coursepartposition{Introduction}%
  \gdef\coursepartindex{0}%
  \gdef\coursepartcolor{courseproblem}}
\newcommand{\coursereperes}{%
  \coursereperemodetrue
  \setcounter{coursechapter}{0}%
  \renewcommand{\coursechapterkind}{Repère}%
  \renewcommand{\thecoursechapter}{\Roman{coursechapter}}}
\newcommand{\coursemain}{%
  \coursereperemodefalse
  \setcounter{coursechapter}{0}%
  \renewcommand{\coursechapterkind}{Chapitre}%
  \renewcommand{\thecoursechapter}{\arabic{coursechapter}}}
\newcommand{\coursepart}[3]{%
  \coursefinishstage
  \coursenavigationvisibletrue
  \stepcounter{coursestep}%
  \gdef\coursepartshort{#2}%
  \gdef\coursepartposition{#2}%
  \xdef\coursepartindex{\arabic{coursestep}}%
  \gdef\coursepartcolor{#3}%
  \setbeamercolor{frametitle}{fg=white,bg=ink}%
  \part{#1}
  \section{#1}
  \begin{frame}[t]{#2}
    \label{course-stage-start-\coursepartindex}%
    \centering\vfill
    {\usebeamerfont{part title}\color{ink}#1\par}
    \vspace{3mm}{\color{#3}\rule{28mm}{1pt}}
    \vfill
  \end{frame}}
\newcommand{\courseunit}[2]{%
  \gdef\coursechapterindicator{}%
  \label{#2}}
\newcommand{\coursechapter}[2]{%
  \ifcoursereperemode
    \gdef\coursechapterindicator{}%
    \label{#2}%
    \section{#1}%
  \else
    \refstepcounter{coursechapter}%
    \gdef\coursechapterindicator{\coursechapterkind\space\thecoursechapter}%
    \label{#2}%
    \section{#1}
    \begin{frame}[plain]
      \vfill
      {\large\sffamily\bfseries\color{\coursepartcolor}%
        \coursechapterkind\space\thecoursechapter\par}
      \vspace{3mm}
      {\usebeamerfont{section title}\color{ink}#1\par}
      \vspace{4mm}{\color{\coursepartcolor}\rule{24mm}{.9pt}}
      \vfill
    \end{frame}%
  \fi}
\renewcommand{\thefigure}{\arabic{figure}}
\renewcommand{\thetable}{\arabic{table}}

% Ouverture pédagogique d'un chapitre : objectifs scannables, métadonnées après.
\NewDocumentCommand{\orientobjective}{m}{\item #1}
\NewDocumentCommand{\orientobjectives}{>{\SplitList{;}}m}{%
  \begin{itemize}
  \setlength{\itemsep}{.45em}
  \ProcessList{#1}{\orientobjective}
  \end{itemize}}
\newcommand{\orientmeta}[2]{%
  \par\addvspace{3pt}\noindent
  \hangindent=26mm\hangafter=1
  \makebox[26mm][l]{\small\bfseries\color{accent!78!ink}#1}%
  {\small\color{ink!68}#2}\par}
\newcommand{\orientcontent}[3]{%
  \orientobjectives{#1}
  \vspace{.35em}
  \orientmeta{Prérequis}{#2}
  \orientmeta{Parcours}{#3}}

\newenvironment{coursecallout}[2]{%
  \par\addvspace{.55em}\begingroup
  \def\coursecalloutcolor{#1}%
  \noindent\color{#1}\rule{\linewidth}{.7pt}\par
  \if\relax\detokenize{#2}\relax\else
    \noindent\colorbox{#1!12}{\parbox{\dimexpr\linewidth-2\fboxsep\relax}{%
      \small\bfseries\color{#1}#2}}\par
  \fi
  \smallskip\leftskip=3mm\rightskip=2mm\raggedright\color{ink!94}\ignorespaces}
  {\par\leftskip=0pt\rightskip=0pt\smallskip
   \noindent\color{\coursecalloutcolor}\rule{\linewidth}{.7pt}\par
   \endgroup\addvspace{.55em}}
\newenvironment{resultat}[1][]{\begin{coursecallout}{accent!82!ink}{#1}}
  {\end{coursecallout}}
\newenvironment{vigilance}[1][]{\begin{coursecallout}{warning!82!ink}{#1}}
  {\end{coursecallout}}
\RenewDocumentEnvironment{lecture}{}{\begin{coursecallout}{coursescore!88!ink}{Lecture}}
  {\end{coursecallout}}
\newenvironment{chiffres}[1][]{\begin{coursecallout}{courseprobability!88!ink}{#1}\small}
  {\end{coursecallout}}
\newenvironment{releve}{%
  \par\begingroup\small\color{ink!68}\leftskip=5mm\rightskip=2mm\raggedright}
  {\par\endgroup}
\newenvironment{rappel}[1][]{\begin{coursecallout}{rulecolor!82!ink}{#1}}
  {\end{coursecallout}}
\newenvironment{approfondir}[1][]{\begin{coursecallout}{courseprobability!82!ink}{#1}}
  {\end{coursecallout}}

\NewDocumentEnvironment{panelenumerate}{O{accent}}
  {\begin{enumerate}\setlength{\itemsep}{.45em}}
  {\end{enumerate}}
\NewDocumentCommand{\panelitem}{O{accent}m}{\item #2}
\NewDocumentEnvironment{panelitemize}{O{accent}}
  {\begin{itemize}\setlength{\itemsep}{.45em}}
  {\end{itemize}}
\NewDocumentCommand{\panelbullet}{O{accent}m}{\item #2}
\NewDocumentCommand{\paneltablerow}{O{accent}mmm}{%
  \cellcolor{#1!10}{\bfseries\color{#1!78!ink}#2} &
  \cellcolor{gridgray!6}{\color{ink!94}#3} &
  \cellcolor{gridgray!6}{\color{ink!65}#4} \\}

\newcounter{exo}
\renewcommand{\theexo}{\arabic{exo}}
\NewDocumentEnvironment{exercice}{O{}}{%
  \refstepcounter{exo}\begin{coursecallout}{rulecolor!82!ink}%
    {Exercice~\theexo\space— #1}}
  {\end{coursecallout}}
\newcounter{tache}
\NewDocumentEnvironment{tache}{O{}}{%
  \refstepcounter{tache}\begin{coursecallout}{accent!82!ink}%
    {Étape~\thetache\space— #1}}
  {\end{coursecallout}}
\newcommand{\tlab}[1]{\par\addvspace{3pt}%
  {\scriptsize\bfseries\color{accent!80!ink}#1}\quad\ignorespaces}
\newcommand{\corrige}[1]{}

\NewDocumentEnvironment{graphique}{m}{}{}
\newcommand{\codefile}[2][firstline=1]{%
  \par\lstinputlisting[style=code,#1]{#2}\par}
\lstnewenvironment{console}
  {\lstset{language=,basicstyle=\ttfamily\scriptsize,
    columns=fullflexible,keepspaces=true}}
  {}

\newcommand{\ann}[3]{%
  \underbrace{\textcolor{#1}{#2}}_{\textcolor{#1}{\substack{#3}}}}
\newcommand{\annup}[3]{%
  \overbrace{\textcolor{#1}{#2}}^{\textcolor{#1}{\substack{#3}}}}
\newcommand{\annl}[1]{\text{\scriptsize #1}}
\newcommand{\R}{\mathbb{R}}
\newcommand{\Pp}{\mathbb{P}}
\newcommand{\x}{\mathbf{x}}
\newcommand{\w}{\mathbf{w}}
\newcommand{\X}{\mathbf{X}}
\newcommand{\y}{\mathbf{y}}
\newcommand{\slidebreak}{\framebreak}

\hypersetup{
  colorlinks=true,
  linkcolor=accent!85!ink,
  urlcolor=accent,
  pdftitle={Régression logistique binaire},
  pdfauthor={Sylvain Maitre},
  pdflang={fr-FR}}

\begin{document}

\begin{frame}[plain]
  \titlepage
\end{frame}

% Fichier généré par scripts/assembler_corps.py.
% Les nombres viennent exclusivement de data/eleves.csv.
% Fichier genere par scripts/construire_cas.py ; ne pas modifier.
% Les donnees brutes, l'imputation et les calculs viennent de data/eleves.csv.

\newcommand{\MicroGryffondorLabel}{Gryffondor}
\newcommand{\MicroSerpentardLabel}{Serpentard}
\newcommand{\MicroNTrain}{20}
\newcommand{\MicroNTest}{8}
\newcommand{\MicroNTotal}{28}
\newcommand{\MicroTrainGryffondor}{12}
\newcommand{\MicroTrainSerpentard}{8}
\newcommand{\MicroMedianPotions}{10}
\newcommand{\MicroMedianVol}{57}
\newcommand{\MicroMeanPotions}{10}
\newcommand{\MicroMeanVol}{50}
\newcommand{\MicroSdPotions}{4}
\newcommand{\MicroSdVol}{20}
\newcommand{\MicroImputedStdPotions}{0}
\newcommand{\MicroImputedStdVol}{0{,}35}
\newcommand{\MicroTrainCovariance}{-64}
\newcommand{\MicroTrainCorrelation}{-\frac{4}{5}}
\newcommand{\MicroTotalStudents}{28}
\newcommand{\MicroWDecimal}{\left(+0{,}603740\,;\,-1{,}029417\,;\,+0{,}514094\right)}
\newcommand{\MicroRawBetaDecimal}{\left(+1{,}892048\,;\,-0{,}257354\,;\,+0{,}025705\right)}
\newcommand{\MicroRawScore}{z=+1{,}892048-0{,}257354\,P+0{,}025705\,V}
\newcommand{\MicroBoundaryRaw}{V=10{,}0120\,P-73{,}6071}
\newcommand{\MicroBoundarySlope}{10.011961}
\newcommand{\MicroBoundaryIntercept}{-73.607139}
\newcommand{\MicroBoundaryInterceptHigh}{-49.524400}
\newcommand{\MicroBoundaryShift}{24{,}083}
\newcommand{\MicroHarryPotions}{-0{,}50}
\newcommand{\MicroHarryVol}{+0{,}50}
\newcommand{\MicroPotionsMin}{4}
\newcommand{\MicroPotionsMax}{17}
\newcommand{\MicroVolMin}{9}
\newcommand{\MicroVolMax}{85}
\newcommand{\MicroTrialWeights}{(0,-2,+1)}
\newcommand{\MicroTrialScore}{z=x_{\mathrm{vol}}-2x_{\mathrm{pot}}}
\newcommand{\MicroTrialRows}{
Harry & $-0{,}50$ & $+0{,}50$ & $+1{,}00$ & $+0{,}50$ & $+1{,}50$ & $0{,}818$ & $0{,}201$ \\
Drago & $+1{,}25$ & $-1{,}05$ & $-2{,}50$ & $-1{,}05$ & $-3{,}55$ & $0{,}028$ & $0{,}028$ \\
Hermione & $+1{,}00$ & $-1{,}10$ & $-2{,}00$ & $-1{,}10$ & $-3{,}10$ & $0{,}043$ & $3{,}144$ \\
Marcus & $-1{,}00$ & $+0{,}75$ & $+2{,}00$ & $+0{,}75$ & $+2{,}75$ & $0{,}940$ & $2{,}812$ \\
Neville & $-0{,}75$ & $+0{,}10$ & $+1{,}50$ & $+0{,}10$ & $+1{,}60$ & $0{,}832$ & $0{,}184$ \\}
\newcommand{\MicroTrialProbabilityRows}{
Harry & \textcolor{coursescore!88!ink}{Gryffondor} & $+1{,}50$ & $0{,}818$ & \textcolor{coursescore!88!ink}{Gryffondor} \\
Drago & \textcolor{warning!88!ink}{Serpentard} & $-3{,}55$ & $0{,}028$ & \textcolor{warning!88!ink}{Serpentard} \\
Hermione & \textcolor{coursescore!88!ink}{Gryffondor} & $-3{,}10$ & $0{,}043$ & \textcolor{warning!88!ink}{Serpentard}\;\textcolor{warning}{$\bullet$} \\
Marcus & \textcolor{warning!88!ink}{Serpentard} & $+2{,}75$ & $0{,}940$ & \textcolor{coursescore!88!ink}{Gryffondor}\;\textcolor{warning}{$\bullet$} \\
Neville & \textcolor{coursescore!88!ink}{Gryffondor} & $+1{,}60$ & $0{,}832$ & \textcolor{coursescore!88!ink}{Gryffondor} \\}
\newcommand{\MicroTrialLoss}{0{,}617793}
\newcommand{\MicroTrialAgree}{15}
\newcommand{\MicroTrialDisagree}{5}
\newcommand{\MicroTrialBestName}{Harry}
\newcommand{\MicroTrialBestScore}{+1{,}50}
\newcommand{\MicroTrialBestLoss}{0{,}201}
\newcommand{\MicroTrialWorstName}{Hermione}
\newcommand{\MicroTrialWorstScore}{-3{,}10}
\newcommand{\MicroTrialWorstLoss}{3{,}144}
\newcommand{\MicroTrialLossRatio}{16}
\newcommand{\MicroResidualRows}{
Harry & $1$ & $-0{,}5$ & $-0{,}50$ & $+0{,}250$ & $-0{,}250$ \\
Drago & $0$ & $+0{,}5$ & $+1{,}25$ & $+0{,}625$ & $-0{,}525$ \\
Hermione & $1$ & $-0{,}5$ & $+1{,}00$ & $-0{,}500$ & $+0{,}550$ \\
Marcus & $0$ & $+0{,}5$ & $-1{,}00$ & $-0{,}500$ & $+0{,}375$ \\
Neville & $1$ & $-0{,}5$ & $-0{,}75$ & $+0{,}375$ & $-0{,}050$ \\}
\newcommand{\MicroGradientZero}{\left(-0{,}100000\,;\,+0{,}275000\,;\,-0{,}257500\right)}
\newcommand{\MicroGradientZeroSumOne}{-2{,}0000}
\newcommand{\MicroGradientZeroSumPotions}{+5{,}5000}
\newcommand{\MicroGradientZeroSumVol}{-5{,}1500}
\newcommand{\MicroFirstStep}{\left(+0{,}100000\,;\,-0{,}275000\,;\,+0{,}257500\right)}
\newcommand{\MicroLossAtZero}{0{,}693147}
\newcommand{\MicroLossAfterStep}{0{,}573791}
\newcommand{\MicroCheapName}{Angelina}
\newcommand{\MicroCheapProbability}{0{,}9407}
\newcommand{\MicroCheapLoss}{0{,}0611}
\newcommand{\MicroBorderName}{Seamus}
\newcommand{\MicroBorderProbability}{0{,}5857}
\newcommand{\MicroBorderLoss}{0{,}5350}
\newcommand{\MicroExpensiveName}{Marcus}
\newcommand{\MicroExpensiveProbability}{0{,}8827}
\newcommand{\MicroExpensiveLoss}{2{,}1434}
\newcommand{\MicroFirstOrderRows}{
$1$ & $\sum_i(p_i-y_i)\,1$ & $2{,}776\times10^{-17}$ \\
$x_{\mathrm{pot}}$ & $\sum_i(p_i-y_i)\,x_{\mathrm{pot}}$ & $7{,}494\times10^{-16}$ \\
$x_{\mathrm{vol}}$ & $\sum_i(p_i-y_i)\,x_{\mathrm{vol}}$ & $5{,}516\times10^{-16}$ \\}
\newcommand{\MicroLossCurveExpr}{(ln(1+exp(-1.375495*x))+ln(1+exp(-1.222830*x))+ln(1+exp(+0.991181*x))+ln(1+exp(-0.579291*x))+ln(1+exp(-0.783672*x))+ln(1+exp(-1.479570*x))+ln(1+exp(-1.427212*x))+ln(1+exp(-2.404297*x))+ln(1+exp(-0.425063*x))+ln(1+exp(-0.346078*x))+ln(1+exp(-2.327491*x))+ln(1+exp(+2.018727*x))+ln(1+exp(+0.682417*x))+ln(1+exp(-2.764470*x))+ln(1+exp(-2.687664*x))+ln(1+exp(-1.041027*x))+ln(1+exp(-1.660425*x))+ln(1+exp(-1.094614*x))+ln(1+exp(+0.963605*x))+ln(1+exp(-2.070444*x)))/20}
\newcommand{\MicroWrongSideNames}{Hermione, Marcus, Colin, Millicent}
\newcommand{\MicroWrongSideCount}{4}
\newcommand{\MicroWrongSideLabels}{
\node[font=\scriptsize\bfseries,text=warning,fill=white,fill opacity=.85,text opacity=1,inner sep=1pt,xshift=0pt,yshift=-8pt] at (axis cs:14,28) {Hermione};
\node[font=\scriptsize\bfseries,text=warning,fill=white,fill opacity=.85,text opacity=1,inner sep=1pt,xshift=0pt,yshift=8pt] at (axis cs:6,65) {Marcus};
\node[font=\scriptsize\bfseries,text=warning,fill=white,fill opacity=.85,text opacity=1,inner sep=1pt,xshift=0pt,yshift=8pt] at (axis cs:13,30) {Colin};
\node[font=\scriptsize\bfseries,text=warning,fill=white,fill opacity=.85,text opacity=1,inner sep=1pt,xshift=0pt,yshift=8pt] at (axis cs:10,64) {Millicent};}
\newcommand{\MicroThresholdBandLabels}{
\node[font=\scriptsize\bfseries,text=warning,fill=white,fill opacity=.85,text opacity=1,inner sep=1pt,xshift=0pt,yshift=8pt] at (axis cs:11,54) {Olivier};
\node[font=\scriptsize\bfseries,text=warning,fill=white,fill opacity=.85,text opacity=1,inner sep=1pt,xshift=0pt,yshift=8pt] at (axis cs:10,41) {Daphné};}
\newcommand{\MicroTrainLoss}{0{,}489080}
\newcommand{\MicroTrainLossPlain}{0.489080}
\newcommand{\MicroTestLoss}{0{,}420942}
\newcommand{\MicroGDAlpha}{1{,}00}
\newcommand{\MicroGDIterations}{690}
\newcommand{\MicroGDLoss}{0{,}489080}
\newcommand{\MicroGDW}{\left(+0{,}603740\,;\,-1{,}029417\,;\,+0{,}514094\right)}
\newcommand{\MicroGradientMaxError}{0{,}00000000000}
\newcommand{\MicroThresholdLow}{0{,}50}
\newcommand{\MicroThresholdHigh}{0{,}65}
\newcommand{\MicroTrainAccuracy}{0{,}8000}
\newcommand{\MicroTrainBaseline}{0{,}6000}
\newcommand{\MicroTestAccuracy}{0{,}7500}
\newcommand{\MicroTestAccuracyTauHigh}{0{,}7500}
\newcommand{\MicroTestPrecision}{0{,}7500}
\newcommand{\MicroTestRecall}{0{,}7500}
\newcommand{\MicroTestSpecificity}{0{,}7500}
\newcommand{\MicroTestFScore}{0{,}7500}
\newcommand{\MicroTestPrecisionTauHigh}{1{,}0000}
\newcommand{\MicroTestRecallTauHigh}{0{,}5000}
\newcommand{\MicroTestWilsonLow}{0{,}409275}
\newcommand{\MicroTestWilsonHigh}{0{,}928521}
\newcommand{\MicroTrainTP}{10}
\newcommand{\MicroTrainFN}{2}
\newcommand{\MicroTrainFP}{2}
\newcommand{\MicroTrainTN}{6}
\newcommand{\MicroTestTP}{3}
\newcommand{\MicroTestFN}{1}
\newcommand{\MicroTestFP}{1}
\newcommand{\MicroTestTN}{3}
\newcommand{\MicroTestTPHigh}{2}
\newcommand{\MicroTestFNHigh}{2}
\newcommand{\MicroTestFPHigh}{0}
\newcommand{\MicroTestTNHigh}{4}
\newcommand{\MicroTestCorrect}{6}
\newcommand{\MicroTrainConfusion}{VP $=10$, FN $=2$, FP $=2$, VN $=6$}
\newcommand{\MicroTestConfusion}{VP $=3$, FN $=1$, FP $=1$, VN $=3$}
\newcommand{\MicroTestConfusionTauHigh}{VP $=2$, FN $=2$, FP $=0$, VN $=4$}
\newcommand{\MicroTrainRows}{
E01 & Harry Potter & \textcolor{coursescore!88!ink}{Gryffondor} & $8$ & $60$ \\
E02 & Drago Malefoy & \textcolor{warning!88!ink}{Serpentard} & $15$ & $29$ \\
E03 & Hermione Granger & \textcolor{coursescore!88!ink}{Gryffondor} & $14$ & $28$ \\
E04 & Pansy Parkinson & \textcolor{warning!88!ink}{Serpentard} & $12$ & $24$ \\
E05 & Ron Weasley & \textcolor{coursescore!88!ink}{Gryffondor} & \textcolor{warning}{---} & $57$ \\
E06 & Gregory Goyle & \textcolor{warning!88!ink}{Serpentard} & $14$ & $9$ \\
E07 & Neville Londubat & \textcolor{coursescore!88!ink}{Gryffondor} & $7$ & $52$ \\
E08 & Ginny Weasley & \textcolor{coursescore!88!ink}{Gryffondor} & $6$ & $80$ \\
E09 & Vincent Crabbe & \textcolor{warning!88!ink}{Serpentard} & $12$ & $30$ \\
E10 & Seamus Finnigan & \textcolor{coursescore!88!ink}{Gryffondor} & $12$ & $60$ \\
E11 & Dean Thomas & \textcolor{coursescore!88!ink}{Gryffondor} & $5$ & $67$ \\
E12 & Marcus Flint & \textcolor{warning!88!ink}{Serpentard} & $6$ & $65$ \\
E13 & Colin Crivey & \textcolor{coursescore!88!ink}{Gryffondor} & $13$ & $30$ \\
E14 & Angelina Johnson & \textcolor{coursescore!88!ink}{Gryffondor} & $5$ & $84$ \\
E15 & Lavande Brown & \textcolor{coursescore!88!ink}{Gryffondor} & $4$ & $71$ \\
E16 & Katie Bell & \textcolor{coursescore!88!ink}{Gryffondor} & $9$ & $57$ \\
E17 & Blaise Zabini & \textcolor{warning!88!ink}{Serpentard} & $17$ & $32$ \\
E18 & Théodore Nott & \textcolor{warning!88!ink}{Serpentard} & $16$ & $44$ \\
E19 & Millicent Bulstrode & \textcolor{warning!88!ink}{Serpentard} & $10$ & $64$ \\
E20 & Parvati Patil & \textcolor{coursescore!88!ink}{Gryffondor} & $5$ & $57$ \\}
\newcommand{\MicroTrainPreparedRows}{
E01 & Harry Potter & \textcolor{coursescore!88!ink}{Gryffondor} & $8$ & $60$ \\
E02 & Drago Malefoy & \textcolor{warning!88!ink}{Serpentard} & $15$ & $29$ \\
E03 & Hermione Granger & \textcolor{coursescore!88!ink}{Gryffondor} & $14$ & $28$ \\
E04 & Pansy Parkinson & \textcolor{warning!88!ink}{Serpentard} & $12$ & $24$ \\
E05 & Ron Weasley & \textcolor{coursescore!88!ink}{Gryffondor} & $10$ & $57$ \\
E06 & Gregory Goyle & \textcolor{warning!88!ink}{Serpentard} & $14$ & $9$ \\
E07 & Neville Londubat & \textcolor{coursescore!88!ink}{Gryffondor} & $7$ & $52$ \\
E08 & Ginny Weasley & \textcolor{coursescore!88!ink}{Gryffondor} & $6$ & $80$ \\
E09 & Vincent Crabbe & \textcolor{warning!88!ink}{Serpentard} & $12$ & $30$ \\
E10 & Seamus Finnigan & \textcolor{coursescore!88!ink}{Gryffondor} & $12$ & $60$ \\
E11 & Dean Thomas & \textcolor{coursescore!88!ink}{Gryffondor} & $5$ & $67$ \\
E12 & Marcus Flint & \textcolor{warning!88!ink}{Serpentard} & $6$ & $65$ \\
E13 & Colin Crivey & \textcolor{coursescore!88!ink}{Gryffondor} & $13$ & $30$ \\
E14 & Angelina Johnson & \textcolor{coursescore!88!ink}{Gryffondor} & $5$ & $84$ \\
E15 & Lavande Brown & \textcolor{coursescore!88!ink}{Gryffondor} & $4$ & $71$ \\
E16 & Katie Bell & \textcolor{coursescore!88!ink}{Gryffondor} & $9$ & $57$ \\
E17 & Blaise Zabini & \textcolor{warning!88!ink}{Serpentard} & $17$ & $32$ \\
E18 & Théodore Nott & \textcolor{warning!88!ink}{Serpentard} & $16$ & $44$ \\
E19 & Millicent Bulstrode & \textcolor{warning!88!ink}{Serpentard} & $10$ & $64$ \\
E20 & Parvati Patil & \textcolor{coursescore!88!ink}{Gryffondor} & $5$ & $57$ \\}
\newcommand{\MicroTestRawRows}{
E21 & Fred Weasley & \textcolor{coursescore!88!ink}{Gryffondor} & $8$ & $85$ \\
E22 & Olivier Dubois & \textcolor{coursescore!88!ink}{Gryffondor} & $11$ & $54$ \\
E23 & George Weasley & \textcolor{coursescore!88!ink}{Gryffondor} & $6$ & \textcolor{warning}{---} \\
E24 & Alicia Spinnet & \textcolor{coursescore!88!ink}{Gryffondor} & $14$ & $45$ \\
E25 & Daphné Greengrass & \textcolor{warning!88!ink}{Serpentard} & $10$ & $41$ \\
E26 & Adrian Pucey & \textcolor{warning!88!ink}{Serpentard} & $14$ & $40$ \\
E27 & Graham Montague & \textcolor{warning!88!ink}{Serpentard} & $16$ & $13$ \\
E28 & Terence Higgs & \textcolor{warning!88!ink}{Serpentard} & $17$ & $24$ \\}
\newcommand{\MicroTestPreparedRows}{
E21 & Fred Weasley & \textcolor{coursescore!88!ink}{Gryffondor} & $8$ & $85$ \\
E22 & Olivier Dubois & \textcolor{coursescore!88!ink}{Gryffondor} & $11$ & $54$ \\
E23 & George Weasley & \textcolor{coursescore!88!ink}{Gryffondor} & $6$ & $57$ \\
E24 & Alicia Spinnet & \textcolor{coursescore!88!ink}{Gryffondor} & $14$ & $45$ \\
E25 & Daphné Greengrass & \textcolor{warning!88!ink}{Serpentard} & $10$ & $41$ \\
E26 & Adrian Pucey & \textcolor{warning!88!ink}{Serpentard} & $14$ & $40$ \\
E27 & Graham Montague & \textcolor{warning!88!ink}{Serpentard} & $16$ & $13$ \\
E28 & Terence Higgs & \textcolor{warning!88!ink}{Serpentard} & $17$ & $24$ \\}
\newcommand{\MicroCompactRows}{
\textcolor{coursescore!88!ink}{Harry} & \textcolor{coursescore!88!ink}{Gryffondor} & $8$ & $60$ \\
\textcolor{warning!88!ink}{Drago} & \textcolor{warning!88!ink}{Serpentard} & $15$ & $29$ \\
\textcolor{coursescore!88!ink}{Hermione} & \textcolor{coursescore!88!ink}{Gryffondor} & $14$ & $28$ \\
\textcolor{warning!88!ink}{Pansy} & \textcolor{warning!88!ink}{Serpentard} & $12$ & $24$ \\
\textcolor{coursescore!88!ink}{Ron} & \textcolor{coursescore!88!ink}{Gryffondor} & \textcolor{warning}{---} & $57$ \\
\textcolor{warning!88!ink}{Gregory} & \textcolor{warning!88!ink}{Serpentard} & $14$ & $9$ \\
\textcolor{coursescore!88!ink}{Neville} & \textcolor{coursescore!88!ink}{Gryffondor} & $7$ & $52$ \\
\textcolor{coursescore!88!ink}{Ginny} & \textcolor{coursescore!88!ink}{Gryffondor} & $6$ & $80$ \\
\textcolor{warning!88!ink}{Vincent} & \textcolor{warning!88!ink}{Serpentard} & $12$ & $30$ \\
\textcolor{coursescore!88!ink}{Seamus} & \textcolor{coursescore!88!ink}{Gryffondor} & $12$ & $60$ \\
\textcolor{coursescore!88!ink}{Dean} & \textcolor{coursescore!88!ink}{Gryffondor} & $5$ & $67$ \\
\textcolor{warning!88!ink}{Marcus} & \textcolor{warning!88!ink}{Serpentard} & $6$ & $65$ \\
\textcolor{coursescore!88!ink}{Colin} & \textcolor{coursescore!88!ink}{Gryffondor} & $13$ & $30$ \\
\textcolor{coursescore!88!ink}{Angelina} & \textcolor{coursescore!88!ink}{Gryffondor} & $5$ & $84$ \\
\textcolor{coursescore!88!ink}{Lavande} & \textcolor{coursescore!88!ink}{Gryffondor} & $4$ & $71$ \\
\textcolor{coursescore!88!ink}{Katie} & \textcolor{coursescore!88!ink}{Gryffondor} & $9$ & $57$ \\
\textcolor{warning!88!ink}{Blaise} & \textcolor{warning!88!ink}{Serpentard} & $17$ & $32$ \\
\textcolor{warning!88!ink}{Théodore} & \textcolor{warning!88!ink}{Serpentard} & $16$ & $44$ \\
\textcolor{warning!88!ink}{Millicent} & \textcolor{warning!88!ink}{Serpentard} & $10$ & $64$ \\
\textcolor{coursescore!88!ink}{Parvati} & \textcolor{coursescore!88!ink}{Gryffondor} & $5$ & $57$ \\
\textcolor{coursescore!88!ink}{Fred} & \textcolor{coursescore!88!ink}{Gryffondor} & $8$ & $85$ \\
\textcolor{coursescore!88!ink}{Olivier} & \textcolor{coursescore!88!ink}{Gryffondor} & $11$ & $54$ \\
\textcolor{coursescore!88!ink}{George} & \textcolor{coursescore!88!ink}{Gryffondor} & $6$ & \textcolor{warning}{---} \\
\textcolor{coursescore!88!ink}{Alicia} & \textcolor{coursescore!88!ink}{Gryffondor} & $14$ & $45$ \\
\textcolor{warning!88!ink}{Daphné} & \textcolor{warning!88!ink}{Serpentard} & $10$ & $41$ \\
\textcolor{warning!88!ink}{Adrian} & \textcolor{warning!88!ink}{Serpentard} & $14$ & $40$ \\
\textcolor{warning!88!ink}{Graham} & \textcolor{warning!88!ink}{Serpentard} & $16$ & $13$ \\
\textcolor{warning!88!ink}{Terence} & \textcolor{warning!88!ink}{Serpentard} & $17$ & $24$ \\}
\newcommand{\MicroSigmoidGryffondorCoords}{(1.375495,0.798267) (-0.991181,0.270679) (0.783672,0.686471) (1.427212,0.806467) (2.404297,0.917154) (0.346078,0.585666) (2.327491,0.911128) (-0.682417,0.335722) (2.764470,0.940725) (2.687664,0.936295) (1.041027,0.739048) (2.070444,0.887997)}
\newcommand{\MicroSigmoidSerpentardCoords}{(-1.222830,0.227439) (-0.579291,0.359096) (-1.479570,0.185492) (-0.425063,0.395306) (2.018727,0.882749) (-1.660425,0.159705) (-1.094614,0.250750) (0.963605,0.723843)}
\newcommand{\MicroDecisionGryffondorCoords}{(0.798267,0.660000) (0.270679,0.660000) (0.686471,0.660000) (0.806467,0.660000) (0.917154,0.660000) (0.585666,0.660000) (0.911128,0.660000) (0.335722,0.660000) (0.940725,0.660000) (0.936295,0.660000) (0.739048,0.660000) (0.887997,0.660000)}
\newcommand{\MicroDecisionSerpentardCoords}{(0.227439,0.340000) (0.359096,0.340000) (0.185492,0.340000) (0.395306,0.340000) (0.882749,0.340000) (0.159705,0.340000) (0.250750,0.340000) (0.723843,0.340000)}
\newcommand{\MicroGDPathCoords}{(0.000000,0.000000) (-0.275000,0.257500) (-0.433705,0.399721) (-0.655142,0.577057) (-0.785845,0.643876) (-0.880651,0.634754) (-0.978382,0.561283) (-1.019962,0.522817) (-1.029091,0.514394) (-1.029417,0.514094) (-1.029417,0.514094)}
\newcommand{\MicroOptimumCoords}{(-1.029417,0.514094)}
\newcommand{\MicroContourOuter}{(-0.411715,0.514094) (-0.344912,0.604210) (-0.269389,0.717743) (-0.187289,0.862914) (-0.109945,1.044951) (-0.067950,1.251853) (-0.109702,1.433808) (-0.255771,1.522328) (-0.456236,1.506872) (-0.647139,1.436995) (-0.804169,1.354732) (-0.928994,1.276884) (-1.029417,1.207198) (-1.112489,1.145088) (-1.183442,1.088921) (-1.246018,1.037015) (-1.302941,0.987850) (-1.356276,0.940064) (-1.407695,0.892371) (-1.458660,0.843463) (-1.510569,0.791887) (-1.564881,0.735890) (-1.623238,0.673207) (-1.687561,0.600740) (-1.760079,0.514094) (-1.843003,0.406983) (-1.936981,0.270913) (-2.035638,0.097303) (-2.111320,-0.110543) (-2.102183,-0.309069) (-1.964064,-0.420553) (-1.748755,-0.423365) (-1.536211,-0.363698) (-1.360725,-0.285754) (-1.223390,-0.209825) (-1.115759,-0.141735) (-1.029417,-0.081716) (-0.957999,-0.028381) (-0.896998,0.019898) (-0.843233,0.064604) (-0.794400,0.107033) (-0.748758,0.148332) (-0.704905,0.189582) (-0.661632,0.231882) (-0.617802,0.276447) (-0.572255,0.324731) (-0.523724,0.378594) (-0.470759,0.440545) (-0.411715,0.514094)}
\newcommand{\MicroContourThird}{(-0.587823,0.514094) (-0.539956,0.578532) (-0.485921,0.659723) (-0.427510,0.763412) (-0.373613,0.892722) (-0.347723,1.037176) (-0.384881,1.158630) (-0.494637,1.211033) (-0.636763,1.194191) (-0.768488,1.144032) (-0.875820,1.087326) (-0.960943,1.034204) (-1.029417,0.986762) (-1.086078,0.944476) (-1.134485,0.906212) (-1.177179,0.870822) (-1.216008,0.837279) (-1.252377,0.804660) (-1.287417,0.772093) (-1.322118,0.738691) (-1.357423,0.703468) (-1.394312,0.665238) (-1.433878,0.622469) (-1.477395,0.573071) (-1.526328,0.514094) (-1.582126,0.441328) (-1.645251,0.349081) (-1.711858,0.231418) (-1.765162,0.089311) (-1.766289,-0.051328) (-1.681387,-0.137876) (-1.537177,-0.147632) (-1.389248,-0.109152) (-1.265227,-0.055202) (-1.167625,-0.001705) (-1.090969,0.046561) (-1.029417,0.089209) (-0.978475,0.127150) (-0.934944,0.161516) (-0.896562,0.193352) (-0.861687,0.223577) (-0.829078,0.253006) (-0.797734,0.282410) (-0.766791,0.312573) (-0.735434,0.344363) (-0.702835,0.378819) (-0.668081,0.417274) (-0.630137,0.461527) (-0.587823,0.514094)}
\newcommand{\MicroContourSecond}{(-0.734748,0.514094) (-0.702699,0.557107) (-0.666544,0.611325) (-0.627633,0.680518) (-0.592398,0.766407) (-0.577520,0.860847) (-0.606467,0.937043) (-0.682415,0.966316) (-0.776404,0.952325) (-0.861744,0.918892) (-0.930790,0.882177) (-0.985452,0.848044) (-1.029417,0.817616) (-1.065806,0.790496) (-1.096900,0.765944) (-1.124326,0.743223) (-1.149266,0.721677) (-1.172618,0.700716) (-1.195106,0.679783) (-1.217363,0.658310) (-1.239989,0.635667) (-1.263603,0.611097) (-1.288896,0.583621) (-1.316668,0.551911) (-1.347834,0.514094) (-1.383296,0.467505) (-1.423366,0.408535) (-1.465816,0.333331) (-1.500854,0.241910) (-1.505171,0.149035) (-1.455590,0.087921) (-1.364715,0.077125) (-1.268261,0.100403) (-1.186266,0.135426) (-1.121419,0.170737) (-1.070404,0.202766) (-1.029417,0.231121) (-0.995486,0.256364) (-0.966486,0.279233) (-0.940911,0.300421) (-0.917667,0.320538) (-0.895927,0.340126) (-0.875023,0.359699) (-0.854377,0.379781) (-0.833446,0.400949) (-0.811672,0.423901) (-0.788445,0.449525) (-0.763067,0.479028) (-0.734748,0.514094)}
\newcommand{\MicroContourInner}{(-0.852463,0.514094) (-0.833149,0.539933) (-0.811365,0.572521) (-0.787997,0.614093) (-0.767174,0.665500) (-0.759433,0.721260) (-0.778825,0.764685) (-0.825617,0.779691) (-0.881589,0.770140) (-0.931649,0.750128) (-0.971941,0.728597) (-1.003798,0.708693) (-1.029417,0.690974) (-1.050625,0.675182) (-1.068749,0.660882) (-1.084735,0.647643) (-1.099271,0.635084) (-1.112879,0.622863) (-1.125978,0.610655) (-1.138937,0.598131) (-1.152102,0.584926) (-1.165831,0.570598) (-1.180522,0.554582) (-1.196632,0.536108) (-1.214685,0.514094) (-1.235195,0.487002) (-1.258353,0.452751) (-1.282963,0.409072) (-1.303727,0.355721) (-1.307783,0.300496) (-1.281143,0.262368) (-1.229101,0.253860) (-1.172272,0.266662) (-1.123389,0.287225) (-1.084570,0.308259) (-1.053993,0.327424) (-1.029417,0.344415) (-1.009071,0.359546) (-0.991680,0.373256) (-0.976341,0.385957) (-0.962399,0.398015) (-0.949355,0.409755) (-0.936810,0.421486) (-0.924415,0.433523) (-0.911843,0.446212) (-0.898758,0.459973) (-0.884790,0.475341) (-0.869518,0.493043) (-0.852463,0.514094)}
\newcommand{\MicroGradientProbe}{(-0.433705,0.399721)}
\newcommand{\MicroGradientProbeX}{-0.433705}
\newcommand{\MicroGradientProbeY}{0.399721}
\newcommand{\MicroGradientTipX}{-0.707653}
\newcommand{\MicroGradientTipY}{0.633285}
\newcommand{\MicroSliceXmin}{-2.191}
\newcommand{\MicroSliceXmax}{0.080}
\newcommand{\MicroSliceYmin}{-0.483}
\newcommand{\MicroSliceYmax}{1.582}
\newcommand{\MicroAllGryffondorCoords}{(8,60) (14,28) (10,57) (7,52) (6,80) (12,60) (5,67) (13,30) (5,84) (4,71) (9,57) (5,57) (8,85) (11,54) (6,57) (14,45)}
\newcommand{\MicroAllSerpentardCoords}{(15,29) (12,24) (14,9) (12,30) (6,65) (17,32) (16,44) (10,64) (10,41) (14,40) (16,13) (17,24)}
\newcommand{\MicroTrainGryffondorCoords}{(8,60) (14,28) (10,57) (7,52) (6,80) (12,60) (5,67) (13,30) (5,84) (4,71) (9,57) (5,57)}
\newcommand{\MicroTrainSerpentardCoords}{(15,29) (12,24) (14,9) (12,30) (6,65) (17,32) (16,44) (10,64)}
\newcommand{\MicroTestGryffondorCoords}{(8,85) (11,54) (6,57) (14,45)}
\newcommand{\MicroTestSerpentardCoords}{(10,41) (14,40) (16,13) (17,24)}
\newcommand{\MicroTrainGryffondorStdCoords}{(-0.5,0.5) (1,-1.1) (0,0.35) (-0.75,0.1) (-1,1.5) (0.5,0.5) (-1.25,0.85) (0.75,-1) (-1.25,1.7) (-1.5,1.05) (-0.25,0.35) (-1.25,0.35)}
\newcommand{\MicroTrainSerpentardStdCoords}{(1.25,-1.05) (0.5,-1.3) (1,-2.05) (0.5,-1) (-1,0.75) (1.75,-0.9) (1.5,-0.3) (0,0.7)}
\newcommand{\MicroMissingCoords}{(10,57) (6,57)}
\newcommand{\MicroAllNameLabels}{
\node[font=\tiny\bfseries,text=coursescore!88!ink,fill=white,fill opacity=.82,text opacity=1,inner sep=.6pt,xshift=-5pt,yshift=9pt] at (axis cs:8,60) {Harry};
\node[font=\tiny\bfseries,text=warning!88!ink,fill=white,fill opacity=.82,text opacity=1,inner sep=.6pt,xshift=-5pt,yshift=-9pt] at (axis cs:15,29) {Drago};
\node[font=\tiny\bfseries,text=coursescore!88!ink,fill=white,fill opacity=.82,text opacity=1,inner sep=.6pt,xshift=-9pt,yshift=-15pt] at (axis cs:14,28) {Hermione};
\node[font=\tiny\bfseries,text=warning!88!ink,fill=white,fill opacity=.82,text opacity=1,inner sep=.6pt,xshift=-13pt,yshift=-3pt] at (axis cs:12,24) {Pansy};
\node[font=\tiny\bfseries,text=coursescore!88!ink,fill=white,fill opacity=.82,text opacity=1,inner sep=.6pt,xshift=10pt,yshift=3pt] at (axis cs:10,57) {Ron};
\node[font=\tiny\bfseries,text=warning!88!ink,fill=white,fill opacity=.82,text opacity=1,inner sep=.6pt,xshift=0pt,yshift=8pt] at (axis cs:14,9) {Gregory};
\node[font=\tiny\bfseries,text=coursescore!88!ink,fill=white,fill opacity=.82,text opacity=1,inner sep=.6pt,xshift=-4pt,yshift=-16pt] at (axis cs:7,52) {Neville};
\node[font=\tiny\bfseries,text=coursescore!88!ink,fill=white,fill opacity=.82,text opacity=1,inner sep=.6pt,xshift=0pt,yshift=-8pt] at (axis cs:6,80) {Ginny};
\node[font=\tiny\bfseries,text=warning!88!ink,fill=white,fill opacity=.82,text opacity=1,inner sep=.6pt,xshift=-17pt,yshift=0pt] at (axis cs:12,30) {Vincent};
\node[font=\tiny\bfseries,text=coursescore!88!ink,fill=white,fill opacity=.82,text opacity=1,inner sep=.6pt,xshift=9pt,yshift=9pt] at (axis cs:12,60) {Seamus};
\node[font=\tiny\bfseries,text=coursescore!88!ink,fill=white,fill opacity=.82,text opacity=1,inner sep=.6pt,xshift=-5pt,yshift=-9pt] at (axis cs:5,67) {Dean};
\node[font=\tiny\bfseries,text=warning!88!ink,fill=white,fill opacity=.82,text opacity=1,inner sep=.6pt,xshift=5pt,yshift=9pt] at (axis cs:6,65) {Marcus};
\node[font=\tiny\bfseries,text=coursescore!88!ink,fill=white,fill opacity=.82,text opacity=1,inner sep=.6pt,xshift=-6pt,yshift=11pt] at (axis cs:13,30) {Colin};
\node[font=\tiny\bfseries,text=coursescore!88!ink,fill=white,fill opacity=.82,text opacity=1,inner sep=.6pt,xshift=0pt,yshift=8pt] at (axis cs:5,84) {Angelina};
\node[font=\tiny\bfseries,text=coursescore!88!ink,fill=white,fill opacity=.82,text opacity=1,inner sep=.6pt,xshift=-3pt,yshift=10pt] at (axis cs:4,71) {Lavande};
\node[font=\tiny\bfseries,text=coursescore!88!ink,fill=white,fill opacity=.82,text opacity=1,inner sep=.6pt,xshift=0pt,yshift=-8pt] at (axis cs:9,57) {Katie};
\node[font=\tiny\bfseries,text=warning!88!ink,fill=white,fill opacity=.82,text opacity=1,inner sep=.6pt,xshift=17pt,yshift=0pt] at (axis cs:17,32) {Blaise};
\node[font=\tiny\bfseries,text=warning!88!ink,fill=white,fill opacity=.82,text opacity=1,inner sep=.6pt,xshift=0pt,yshift=-8pt] at (axis cs:16,44) {Théodore};
\node[font=\tiny\bfseries,text=warning!88!ink,fill=white,fill opacity=.82,text opacity=1,inner sep=.6pt,xshift=0pt,yshift=8pt] at (axis cs:10,64) {Millicent};
\node[font=\tiny\bfseries,text=coursescore!88!ink,fill=white,fill opacity=.82,text opacity=1,inner sep=.6pt,xshift=0pt,yshift=-8pt] at (axis cs:5,57) {Parvati};
\node[font=\tiny\bfseries,text=coursescore!88!ink,fill=white,fill opacity=.82,text opacity=1,inner sep=.6pt,xshift=0pt,yshift=8pt] at (axis cs:8,85) {Fred};
\node[font=\tiny\bfseries,text=coursescore!88!ink,fill=white,fill opacity=.82,text opacity=1,inner sep=.6pt,xshift=3pt,yshift=-10pt] at (axis cs:11,54) {Olivier};
\node[font=\tiny\bfseries,text=coursescore!88!ink,fill=white,fill opacity=.82,text opacity=1,inner sep=.6pt,xshift=-4pt,yshift=-16pt] at (axis cs:6,57) {George};
\node[font=\tiny\bfseries,text=coursescore!88!ink,fill=white,fill opacity=.82,text opacity=1,inner sep=.6pt,xshift=0pt,yshift=8pt] at (axis cs:14,45) {Alicia};
\node[font=\tiny\bfseries,text=warning!88!ink,fill=white,fill opacity=.82,text opacity=1,inner sep=.6pt,xshift=0pt,yshift=-8pt] at (axis cs:10,41) {Daphné};
\node[font=\tiny\bfseries,text=warning!88!ink,fill=white,fill opacity=.82,text opacity=1,inner sep=.6pt,xshift=5pt,yshift=-9pt] at (axis cs:14,40) {Adrian};
\node[font=\tiny\bfseries,text=warning!88!ink,fill=white,fill opacity=.82,text opacity=1,inner sep=.6pt,xshift=0pt,yshift=-8pt] at (axis cs:16,13) {Graham};
\node[font=\tiny\bfseries,text=warning!88!ink,fill=white,fill opacity=.82,text opacity=1,inner sep=.6pt,xshift=0pt,yshift=-8pt] at (axis cs:17,24) {Terence};}
\newcommand{\MicroTrainAtypicalCoords}{(14,28) (12,60) (6,65) (13,30) (10,64)}
\newcommand{\MicroTestAtypicalCoords}{(11,54) (14,45) (10,41)}
\newcommand{\MicroGDTableRows}{
$0$ & $0{,}693147$ & $+0{,}000000$ & $+0{,}000000$ & $+0{,}000000$ & $0{,}389784$ \\
$1$ & $0{,}573791$ & $+0{,}100000$ & $-0{,}275000$ & $+0{,}257500$ & $0{,}226364$ \\
$2$ & $0{,}531952$ & $+0{,}176332$ & $-0{,}433705$ & $+0{,}399721$ & $0{,}145894$ \\
$5$ & $0{,}499775$ & $+0{,}331248$ & $-0{,}655142$ & $+0{,}577057$ & $0{,}058724$ \\
$10$ & $0{,}491770$ & $+0{,}469108$ & $-0{,}785845$ & $+0{,}643876$ & $0{,}022873$ \\
$20$ & $0{,}489764$ & $+0{,}567160$ & $-0{,}880651$ & $+0{,}634754$ & $0{,}007756$ \\
$50$ & $0{,}489160$ & $+0{,}600309$ & $-0{,}978382$ & $+0{,}561283$ & $0{,}002314$ \\
$100$ & $0{,}489082$ & $+0{,}603157$ & $-1{,}019962$ & $+0{,}522817$ & $4{,}268\times10^{-4}$ \\
$200$ & $0{,}489080$ & $+0{,}603719$ & $-1{,}029091$ & $+0{,}514394$ & $1{,}468\times10^{-5}$ \\
$400$ & $0{,}489080$ & $+0{,}603740$ & $-1{,}029417$ & $+0{,}514094$ & $1{,}742\times10^{-8}$ \\
$690$ & $0{,}489080$ & $+0{,}603740$ & $-1{,}029417$ & $+0{,}514094$ & $9{,}967\times10^{-13}$ \\}
\newcommand{\MicroGradientCheckRows}{
$0$ & $-0{,}054724771$ & $-0{,}054724771$ & $0{,}00000000000$ \\
$1$ & $+0{,}124493338$ & $+0{,}124493338$ & $0{,}00000000000$ \\
$2$ & $-0{,}111803507$ & $-0{,}111803507$ & $0{,}00000000000$ \\}
\newcommand{\MicroMetricsRows}{
apprentissage & $0{,}50$ & $10$ & $2$ & $2$ & $6$ & $0{,}8000$ & $0{,}8333$ & $0{,}8333$ & $0{,}7500$ & $0{,}8333$ & $0{,}7917$ \\
apprentissage & $0{,}65$ & $9$ & $3$ & $2$ & $6$ & $0{,}7500$ & $0{,}8182$ & $0{,}7500$ & $0{,}7500$ & $0{,}7826$ & $0{,}7500$ \\
test & $0{,}50$ & $3$ & $1$ & $1$ & $3$ & $0{,}7500$ & $0{,}7500$ & $0{,}7500$ & $0{,}7500$ & $0{,}7500$ & $0{,}7500$ \\
test & $0{,}65$ & $2$ & $2$ & $0$ & $4$ & $0{,}7500$ & $1{,}0000$ & $0{,}5000$ & $1{,}0000$ & $0{,}6667$ & $0{,}7500$ \\}
\newcommand{\MicroTestSummaryRows}{
$0{,}50$ & $0{,}7500$ & $0{,}7500$ & $0{,}7500$ & $0{,}7500$ & $0{,}7500$ \\
$0{,}65$ & $0{,}7500$ & $1{,}0000$ & $0{,}5000$ & $1{,}0000$ & $0{,}6667$ \\}
\newcommand{\MicroTrainCalcRows}{
E01 & Harry & \textcolor{coursescore!88!ink}{Gryff.} & $+1{,}375495$ & $0{,}798267$ & \textcolor{coursescore!88!ink}{Gryff.} & \textcolor{coursescore!88!ink}{Gryff.} \\
E02 & Drago & \textcolor{warning!88!ink}{Serp.} & $-1{,}222830$ & $0{,}227439$ & \textcolor{warning!88!ink}{Serp.} & \textcolor{warning!88!ink}{Serp.} \\
E03 & Hermione & \textcolor{coursescore!88!ink}{Gryff.} & $-0{,}991181$ & $0{,}270679$ & \textcolor{warning!88!ink}{Serp.} & \textcolor{warning!88!ink}{Serp.} \\
E04 & Pansy & \textcolor{warning!88!ink}{Serp.} & $-0{,}579291$ & $0{,}359096$ & \textcolor{warning!88!ink}{Serp.} & \textcolor{warning!88!ink}{Serp.} \\
E05 & Ron & \textcolor{coursescore!88!ink}{Gryff.} & $+0{,}783672$ & $0{,}686471$ & \textcolor{coursescore!88!ink}{Gryff.} & \textcolor{coursescore!88!ink}{Gryff.} \\
E06 & Gregory & \textcolor{warning!88!ink}{Serp.} & $-1{,}479570$ & $0{,}185492$ & \textcolor{warning!88!ink}{Serp.} & \textcolor{warning!88!ink}{Serp.} \\
E07 & Neville & \textcolor{coursescore!88!ink}{Gryff.} & $+1{,}427212$ & $0{,}806467$ & \textcolor{coursescore!88!ink}{Gryff.} & \textcolor{coursescore!88!ink}{Gryff.} \\
E08 & Ginny & \textcolor{coursescore!88!ink}{Gryff.} & $+2{,}404297$ & $0{,}917154$ & \textcolor{coursescore!88!ink}{Gryff.} & \textcolor{coursescore!88!ink}{Gryff.} \\
E09 & Vincent & \textcolor{warning!88!ink}{Serp.} & $-0{,}425063$ & $0{,}395306$ & \textcolor{warning!88!ink}{Serp.} & \textcolor{warning!88!ink}{Serp.} \\
E10 & Seamus & \textcolor{coursescore!88!ink}{Gryff.} & $+0{,}346078$ & $0{,}585666$ & \textcolor{coursescore!88!ink}{Gryff.} & \textcolor{warning!88!ink}{Serp.} \\
E11 & Dean & \textcolor{coursescore!88!ink}{Gryff.} & $+2{,}327491$ & $0{,}911128$ & \textcolor{coursescore!88!ink}{Gryff.} & \textcolor{coursescore!88!ink}{Gryff.} \\
E12 & Marcus & \textcolor{warning!88!ink}{Serp.} & $+2{,}018727$ & $0{,}882749$ & \textcolor{coursescore!88!ink}{Gryff.} & \textcolor{coursescore!88!ink}{Gryff.} \\
E13 & Colin & \textcolor{coursescore!88!ink}{Gryff.} & $-0{,}682417$ & $0{,}335722$ & \textcolor{warning!88!ink}{Serp.} & \textcolor{warning!88!ink}{Serp.} \\
E14 & Angelina & \textcolor{coursescore!88!ink}{Gryff.} & $+2{,}764470$ & $0{,}940725$ & \textcolor{coursescore!88!ink}{Gryff.} & \textcolor{coursescore!88!ink}{Gryff.} \\
E15 & Lavande & \textcolor{coursescore!88!ink}{Gryff.} & $+2{,}687664$ & $0{,}936295$ & \textcolor{coursescore!88!ink}{Gryff.} & \textcolor{coursescore!88!ink}{Gryff.} \\
E16 & Katie & \textcolor{coursescore!88!ink}{Gryff.} & $+1{,}041027$ & $0{,}739048$ & \textcolor{coursescore!88!ink}{Gryff.} & \textcolor{coursescore!88!ink}{Gryff.} \\
E17 & Blaise & \textcolor{warning!88!ink}{Serp.} & $-1{,}660425$ & $0{,}159705$ & \textcolor{warning!88!ink}{Serp.} & \textcolor{warning!88!ink}{Serp.} \\
E18 & Théodore & \textcolor{warning!88!ink}{Serp.} & $-1{,}094614$ & $0{,}250750$ & \textcolor{warning!88!ink}{Serp.} & \textcolor{warning!88!ink}{Serp.} \\
E19 & Millicent & \textcolor{warning!88!ink}{Serp.} & $+0{,}963605$ & $0{,}723843$ & \textcolor{coursescore!88!ink}{Gryff.} & \textcolor{coursescore!88!ink}{Gryff.} \\
E20 & Parvati & \textcolor{coursescore!88!ink}{Gryff.} & $+2{,}070444$ & $0{,}887997$ & \textcolor{coursescore!88!ink}{Gryff.} & \textcolor{coursescore!88!ink}{Gryff.} \\}
\newcommand{\MicroTestCalcRows}{
E21 & Fred & \textcolor{coursescore!88!ink}{Gryff.} & $+2{,}018112$ & $0{,}882686$ & \textcolor{coursescore!88!ink}{Gryff.} & \textcolor{coursescore!88!ink}{Gryff.} \\
E22 & Olivier & \textcolor{coursescore!88!ink}{Gryff.} & $+0{,}449204$ & $0{,}610450$ & \textcolor{coursescore!88!ink}{Gryff.} & \textcolor{warning!88!ink}{Serp.} \\
E23 & George & \textcolor{coursescore!88!ink}{Gryff.} & $+1{,}813090$ & $0{,}859735$ & \textcolor{coursescore!88!ink}{Gryff.} & \textcolor{coursescore!88!ink}{Gryff.} \\
E24 & Alicia & \textcolor{coursescore!88!ink}{Gryff.} & $-0{,}554201$ & $0{,}364890$ & \textcolor{warning!88!ink}{Serp.} & \textcolor{warning!88!ink}{Serp.} \\
E25 & Daphné & \textcolor{warning!88!ink}{Serp.} & $+0{,}372397$ & $0{,}592038$ & \textcolor{coursescore!88!ink}{Gryff.} & \textcolor{warning!88!ink}{Serp.} \\
E26 & Adrian & \textcolor{warning!88!ink}{Serp.} & $-0{,}682724$ & $0{,}335654$ & \textcolor{warning!88!ink}{Serp.} & \textcolor{warning!88!ink}{Serp.} \\
E27 & Graham & \textcolor{warning!88!ink}{Serp.} & $-1{,}891459$ & $0{,}131078$ & \textcolor{warning!88!ink}{Serp.} & \textcolor{warning!88!ink}{Serp.} \\
E28 & Terence & \textcolor{warning!88!ink}{Serp.} & $-1{,}866062$ & $0{,}133998$ & \textcolor{warning!88!ink}{Serp.} & \textcolor{warning!88!ink}{Serp.} \\}
\newcommand{\MicroProbabilityGeorge}{0{,}859735}
\newcommand{\MicroProbabilityDaphne}{0{,}592038}
\newcommand{\MicroProbabilityAlicia}{0{,}364890}
\newcommand{\MicroProbabilityFred}{0{,}882686}
\newcommand{\MicroProbabilityOlivier}{0{,}610450}

% Contenu éditorial du cas Gryffondor--Serpentard.
% Toutes les valeurs numériques affichées sont fournies par le générateur.

% Maisons et matières portent partout la teinte qui les désigne dans les
% figures. Les noms sont écrits en entier chaque fois que la place le permet ;
% les tableaux les plus serrés emploient les formes abrégées de même couleur.
\newcommand{\GryffondorName}{\textcolor{coursescore!88!ink}{\MicroGryffondorLabel}}
\newcommand{\SerpentardName}{\textcolor{warning!88!ink}{\MicroSerpentardLabel}}
\newcommand{\PotionsName}{\textcolor{coursepotions!92!ink}{Potions}}
\newcommand{\VolName}{\textcolor{coursevol!92!ink}{Vol}}
\newcommand{\microheader}{\rowcolor{ink!92}\color{white}\bfseries}

% Les deux notes vivent sur des échelles différentes : sans « scale only axis »,
% la boîte s'étire et les deux groupes d'élèves se lisent comme des traînées
% obliques. Le rapport 1,25 entre largeur et hauteur rend un disque de notes
% rond à l'écran.
\pgfplotsset{
  rawcloud/.style={
    scale only axis,width=7cm,height=5.6cm,
    xmin=0,xmax=20,ymin=0,ymax=100,
    xtick={0,5,10,15,20},ytick={0,20,40,60,80,100},
    axis lines=box,axis line style={ink!42},
    grid=major,grid style={gridgray!34},
    xlabel={\PotionsName{} (sur 20)},ylabel={\VolName{} (sur 100)},
    tick label style={font=\scriptsize},label style={font=\scriptsize},
    legend style={at={(.5,1.02)},anchor=south,draw=none,
      font=\scriptsize,legend columns=3,
      /tikz/every even column/.append style={column sep=7pt}}},
  smallcloud/.style={
    rawcloud,width=4.5cm,height=3.6cm,
    xtick={0,10,20},ytick={0,50,100},
    xlabel={\PotionsName},ylabel={\VolName},
    tick label style={font=\tiny}},
  stdcloud/.style={
    scale only axis,width=4.5cm,height=3.6cm,
    xmin=-2.6,xmax=2.6,ymin=-2.6,ymax=2.6,
    xtick={-2,-1,0,1,2},ytick={-2,-1,0,1,2},
    axis lines=box,axis line style={ink!42},
    grid=major,grid style={gridgray!34},
    xlabel={$x_{\mathrm{pot}}$},ylabel={$x_{\mathrm{vol}}$},
    tick label style={font=\tiny},label style={font=\scriptsize}}}

% Vignettes de la planche d'ensemble. Chaque carte est une image autonome, posée
% par un nœud : son contenu est donc centré par construction, jamais par des
% coordonnées ajustées à la main. L'ensemble est enveloppé dans \hyperref, ce qui
% rend toute la surface de la carte cliquable vers la section correspondante.
\newlength{\chainw}\setlength{\chainw}{3.06cm}
\newlength{\chainh}\setlength{\chainh}{3.6cm}
\pgfplotsset{
  chainaxis/.style={
    scale only axis,width=2.55cm,height=1.7cm,
    axis lines=box,axis line style={ink!30,line width=.4pt},
    xtick=\empty,ytick=\empty,
    clip mode=individual,
    every axis plot/.append style={line width=.7pt}}}

% Lignes de niveau de la perte dans le plan des deux poids, calculées par le
% générateur. Elles servent de fond commun aux vignettes gradient et mise à
% jour : l'une y lit une direction, l'autre y suit une suite de pas.
\newcommand{\microrelief}{%
  \addplot[courseupdate!22!white,line width=.5pt] coordinates {\MicroContourOuter};
  \addplot[courseupdate!34!white,line width=.5pt] coordinates {\MicroContourThird};
  \addplot[courseupdate!48!white,line width=.5pt] coordinates {\MicroContourSecond};
  \addplot[courseupdate!62!white,line width=.5pt] coordinates {\MicroContourInner};}

\newcommand{\chaincard}[4]{%
  \hyperref[#1]{%
    \begin{tikzpicture}
      \path[draw=ink!24,fill=white,rounded corners=1.6pt,line width=.5pt]
        (0,0) rectangle (\chainw,\chainh);
      \path[fill=#2!16,rounded corners=1.6pt]
        (0,\chainh-.51cm) rectangle (\chainw,\chainh);
      \node[font=\fontsize{6.3}{7}\selectfont\bfseries,text=ink!92,anchor=west]
        at (.16,\chainh-.255cm) {#3};
      \node[inner sep=0pt] at (.5\chainw,.5\chainh-.255cm) {#4};
    \end{tikzpicture}}}

% Frontière ajustée et frontière décalée par le seuil haut, toutes deux issues
% des coefficients estimés : aucune constante n'est écrite ici à la main.
\newcommand{\microboundary}[1]{%
  \addplot[#1,domain=0:20,samples=2]
    {\MicroBoundarySlope*x+(\MicroBoundaryIntercept)}}
\newcommand{\microboundaryhigh}[1]{%
  \addplot[#1,domain=0:20,samples=2]
    {\MicroBoundarySlope*x+(\MicroBoundaryInterceptHigh)}}

\courseintro

\begin{frame}[t]{Formulation du problème de classification}
  \vspace{4mm}
  \begin{columns}[T,onlytextwidth]
    \column{.47\textwidth}
      \begin{coursecallout}{coursedata!88!ink}{Observations disponibles}
        Pour chaque élève :
        \begin{itemize}
          \item une note de \PotionsName{} sur 20 ;
          \item une note de \VolName{} sur 100.
        \end{itemize}
      \end{coursecallout}
    \column{.47\textwidth}
      \begin{coursecallout}{coursedecision!88!ink}{Maison observée}
        Deux maisons sont représentées :
        \begin{itemize}
          \item \GryffondorName ;
          \item \SerpentardName.
        \end{itemize}
      \end{coursecallout}
  \end{columns}
  \vfill
  \centering
  {\large\bfseries Construire une règle qui associe les deux notes à une maison.}
\end{frame}

% Les titres de planche passent par \MakeUppercase : aucune commande colorée ne
% peut y figurer, le nom de la couleur serait mis en capitales lui aussi.
\begin{frame}[t]{Distribution des notes de Potions et de Vol}
  \begin{columns}[c,onlytextwidth]
    \column{.34\textwidth}
      \scriptsize
      \renewcommand{\arraystretch}{.66}
      \setlength{\tabcolsep}{3pt}
      \begin{tabular}{@{}l l r r@{}}
        \toprule
        élève & maison & \PotionsName & \VolName\\
        \midrule
        \MicroCompactRows
        \bottomrule
      \end{tabular}
    \column{.64\textwidth}
      \centering
      \begin{tikzpicture}
        \begin{axis}[rawcloud,width=7.15cm,height=5.72cm]
          \addplot[only marks,mark=square*,mark size=2.6pt,coursescore!88!ink]
            coordinates {\MicroAllGryffondorCoords};
          \addlegendentry{\GryffondorName}
          \addplot[only marks,mark=*,mark size=2.6pt,warning!88!ink]
            coordinates {\MicroAllSerpentardCoords};
          \addlegendentry{\SerpentardName}
          \addplot[only marks,mark=o,mark size=4.6pt,very thick,ink!72]
            coordinates {\MicroMissingCoords};
          \addlegendentry{note imputée}
          \MicroAllNameLabels
        \end{axis}
      \end{tikzpicture}
  \end{columns}
\end{frame}

\begin{frame}[t]{Chaîne de construction du classifieur}
\vfill
\centering
\begin{adjustbox}{max totalsize={\textwidth}{.83\textheight},center}
\begin{tikzpicture}[
  flow/.style={-{Latex[length=1.8mm]},line width=.75pt,ink!62},
  slot/.style={inner sep=0pt,outer sep=0pt}]

  \node[slot] (donnees) at (0,4.8) {\chaincard{course-stage-start-1}{coursedata}{DONNÉES}{%
    \begin{tikzpicture}
      \begin{axis}[chainaxis,xmin=0,xmax=20,ymin=0,ymax=100,
        xlabel={\PotionsName},ylabel={\VolName},
        label style={font=\tiny,inner sep=1.5pt}]
        \addplot[only marks,mark=square*,mark size=.9pt,coursescore!88!ink]
          coordinates {\MicroAllGryffondorCoords};
        \addplot[only marks,mark=*,mark size=.9pt,warning!88!ink]
          coordinates {\MicroAllSerpentardCoords};
      \end{axis}
    \end{tikzpicture}}};

  \node[slot] (score) at (3.78,4.8) {\chaincard{course-stage-start-2}{coursescore}{SCORE LINÉAIRE}{%
    \begin{tikzpicture}
      \begin{axis}[chainaxis,xmin=0,xmax=20,ymin=0,ymax=100]
        \addplot[only marks,mark=square*,mark size=.9pt,coursescore!88!ink]
          coordinates {\MicroTrainGryffondorCoords};
        \addplot[only marks,mark=*,mark size=.9pt,warning!88!ink]
          coordinates {\MicroTrainSerpentardCoords};
        \microboundary{ink!72,line width=.8pt};
      \end{axis}
    \end{tikzpicture}}};

  \node[slot] (sigmoide) at (7.56,4.8) {\chaincard{course-stage-start-3}{courseprobability}{SIGMOÏDE}{%
    \begin{tikzpicture}
      \begin{axis}[chainaxis,xmin=-3.2,xmax=3.2,ymin=-.06,ymax=1.06]
        \addplot[courseprobability!88!ink,domain=-3.2:3.2,samples=70]
          {1/(1+exp(-x))};
        \addplot[only marks,mark=square*,mark size=.9pt,coursescore!88!ink]
          coordinates {\MicroSigmoidGryffondorCoords};
        \addplot[only marks,mark=*,mark size=.9pt,warning!88!ink]
          coordinates {\MicroSigmoidSerpentardCoords};
      \end{axis}
    \end{tikzpicture}}};

  \node[slot] (decision) at (11.34,4.8) {\chaincard{course-stage-start-7}{coursedecision}{DÉCISION}{%
    \begin{tikzpicture}
      \begin{axis}[chainaxis,height=.95cm,xmin=0,xmax=1,ymin=0,ymax=1]
        \fill[warning!8] (axis cs:0,0) rectangle (axis cs:.5,1);
        \fill[coursescore!8] (axis cs:.5,0) rectangle (axis cs:1,1);
        \draw[coursedecision!88!ink,dashed,line width=.6pt]
          (axis cs:.5,0) -- (axis cs:.5,1);
        \addplot[only marks,mark=square*,mark size=.9pt,coursescore!88!ink]
          coordinates {\MicroDecisionGryffondorCoords};
        \addplot[only marks,mark=*,mark size=.9pt,warning!88!ink]
          coordinates {\MicroDecisionSerpentardCoords};
        \node[font=\tiny,text=warning!88!ink,anchor=north west,yshift=-1pt]
          at (axis description cs:0,0) {\MicroSerpentardLabel};
        \node[font=\tiny,text=coursescore!88!ink,anchor=north east,yshift=-1pt]
          at (axis description cs:1,0) {\MicroGryffondorLabel};
        \node[font=\tiny,text=coursedecision!88!ink,anchor=north,
          inner sep=1.5pt] at (axis description cs:.5,1) {$\tau$};
      \end{axis}
    \end{tikzpicture}}};

  \node[slot] (maj) at (0,0) {\chaincard{course-stage-start-6}{courseupdate}{MISE À JOUR}{%
    \begin{tikzpicture}
      \begin{axis}[chainaxis,xmin=\MicroSliceXmin,xmax=\MicroSliceXmax,
        ymin=\MicroSliceYmin,ymax=\MicroSliceYmax,
        xlabel={$w_{\mathrm{pot}}$},ylabel={$w_{\mathrm{vol}}$},
        label style={font=\tiny,inner sep=1.5pt}]
        \microrelief
        \addplot[courseupdate!85!ink,line width=.8pt,mark=*,mark size=.75pt,
          mark options={fill=courseupdate!85!ink,draw=none}]
          coordinates {\MicroGDPathCoords};
        \addplot[only marks,mark=star,mark size=2pt,line width=.7pt,ink!75]
          coordinates {\MicroOptimumCoords};
      \end{axis}
    \end{tikzpicture}}};

  \node[slot] (gradient) at (3.78,0) {\chaincard{course-stage-start-5}{coursegradient}{GRADIENT}{%
    \begin{tikzpicture}
      \begin{axis}[chainaxis,xmin=\MicroSliceXmin,xmax=\MicroSliceXmax,
        ymin=\MicroSliceYmin,ymax=\MicroSliceYmax,
        xlabel={$w_{\mathrm{pot}}$},ylabel={$w_{\mathrm{vol}}$},
        label style={font=\tiny,inner sep=1.5pt}]
        \microrelief
        \draw[-{Latex[length=1.7mm]},coursegradient!92!ink,line width=1.1pt]
          (axis cs:\MicroGradientProbeX,\MicroGradientProbeY)
          -- (axis cs:\MicroGradientTipX,\MicroGradientTipY);
        \addplot[only marks,mark=*,mark size=1.2pt,coursegradient!92!ink]
          coordinates {\MicroGradientProbe};
        \node[font=\tiny,text=coursegradient!92!ink,anchor=south west,inner sep=1pt]
          at (axis cs:\MicroGradientTipX,\MicroGradientTipY) {$-\nabla J$};
      \end{axis}
    \end{tikzpicture}}};

  \node[slot] (perte) at (7.56,0) {\chaincard{course-stage-start-4}{courseloss}{PERTE LOGISTIQUE}{%
    \begin{tikzpicture}
      \begin{axis}[chainaxis,xmin=-3,xmax=3,ymin=0,ymax=3.2,
        xlabel={$z$},label style={font=\tiny,inner sep=1.5pt}]
        \addplot[coursescore!88!ink,domain=-3:3,samples=60] {ln(1+exp(-x))};
        \addplot[warning!88!ink,domain=-3:3,samples=60] {ln(1+exp(x))};
        \node[font=\fontsize{5}{5.5}\selectfont,text=coursescore!88!ink,
          anchor=north west,inner sep=2pt]
          at (axis description cs:0,1) {\MicroGryffondorLabel};
        \node[font=\fontsize{5}{5.5}\selectfont,text=warning!88!ink,
          anchor=north east,inner sep=2pt]
          at (axis description cs:1,1) {\MicroSerpentardLabel};
      \end{axis}
    \end{tikzpicture}}};

  \node[slot] (evaluation) at (11.34,0)
    {\chaincard{course-stage-start-8}{courseevaluation}{ÉVALUATION}{%
    \begin{tikzpicture}[x=1.02cm,y=.58cm,
      cell/.style={draw=white,line width=.7pt,minimum width=1cm,
        minimum height=.56cm,font=\tiny\bfseries,inner sep=0pt},
      tag/.style={font=\tiny,text=ink!62}]
      \node[cell,fill=rulecolor!18,text=rulecolor!75!ink] at (0,1) {VN $=\MicroTrainTN$};
      \node[cell,fill=warning!16,text=warning!85!ink]     at (1,1) {FP $=\MicroTrainFP$};
      \node[cell,fill=warning!16,text=warning!85!ink]     at (0,0) {FN $=\MicroTrainFN$};
      \node[cell,fill=rulecolor!18,text=rulecolor!75!ink] at (1,0) {VP $=\MicroTrainTP$};
      \node[tag,anchor=south] at (.5,1.62) {maison prédite};
      \node[tag,anchor=south,rotate=90] at (-.62,.5) {maison observée};
    \end{tikzpicture}}};

  \draw[flow] (donnees) -- (score);
  \draw[flow] (score) -- (sigmoide);
  \draw[flow] (sigmoide) -- (decision);
  \draw[flow] (sigmoide) -- (perte);
  \draw[flow] (perte) -- (gradient);
  \draw[flow] (gradient) -- (maj);
  \draw[flow] (decision) -- (evaluation);
  \draw[flow] (maj.north) -- ++(0,.55) -| (score.south);
\end{tikzpicture}
\end{adjustbox}
\vfill
\end{frame}

% ---------------------------------------------------------------------------
\coursepart{Définir les données et le protocole}{Données}{coursedata}

\begin{frame}[t]{Séparation entre estimation et évaluation}
  \vspace{5mm}
  \centering
  \begin{tikzpicture}[node distance=9mm and 12mm,
    box/.style={draw=ink!28,rounded corners=2pt,minimum height=15mm,
      minimum width=31mm,align=center,font=\small,inner xsep=4mm},
    arr/.style={-{Latex[length=2mm]},thick,ink!58}]
    \node[box,fill=coursedata!13] (all) {24 élèves};
    \node[box,fill=coursedata!22,right=of all,yshift=10mm] (train)
      {\MicroNTrain{} élèves\\estimation};
    \node[box,fill=courseevaluation!18,right=of all,yshift=-10mm] (test)
      {\MicroNTest{} élèves\\évaluation};
    \node[box,fill=courseupdate!15,right=of train] (fit)
      {médianes, $\mu$, $s$\\et coefficients};
    \node[box,fill=courseevaluation!24,right=of test] (eval)
      {mesure\\des erreurs};
    \draw[arr] (all) -- (train);
    \draw[arr] (all) -- (test);
    \draw[arr] (train) -- (fit);
    \draw[arr] (fit.south) |- (eval.north);
    \draw[arr] (test) -- (eval);
  \end{tikzpicture}
  \vfill
  Le groupe d'estimation constitue le \emph{jeu d'apprentissage} : lui seul fixe
  les nombres du modèle. Le groupe d'évaluation les reçoit sans les modifier.
\end{frame}

\begin{frame}[t]{Groupe d'estimation : \MicroNTrain{} élèves}
  \centering
  \tiny
  \renewcommand{\arraystretch}{.9}
  \begin{tabularx}{\textwidth}{c >{\raggedright\arraybackslash}X l c c}
    \microheader ID & \color{white}\bfseries Élève &
      \color{white}\bfseries Maison & \color{white}\bfseries Potions &
      \color{white}\bfseries Vol \\
    \MicroTrainRows
  \end{tabularx}
  \par\vspace{2mm}
  \begin{vigilance}[Une note manque ici même]
    La note de \PotionsName{} de Ron est absente. Le trou se trouve donc là où sont
    calculées les statistiques : sa valeur de remplacement participera aux
    médianes, aux moyennes et aux coefficients.
  \end{vigilance}
\end{frame}

\begin{frame}[t]{Groupe d'évaluation : huit élèves}
  \centering
  \small
  \renewcommand{\arraystretch}{1.04}
  \begin{tabularx}{\textwidth}{c >{\raggedright\arraybackslash}X l c c}
    \microheader ID & \color{white}\bfseries Élève &
      \color{white}\bfseries Maison & \color{white}\bfseries Potions &
      \color{white}\bfseries Vol \\
    \MicroTestRawRows
  \end{tabularx}
  \vfill
  \begin{vigilance}[Une note manque aussi à l'évaluation]
    La note de \VolName{} de George est absente. Une case vide ne peut pas être
    standardisée : elle doit d'abord recevoir une valeur numérique, et cette
    valeur vient du groupe d'estimation.
  \end{vigilance}
\end{frame}

\begin{frame}[t]{Compléter, puis standardiser une note manquante}
  \vspace{3mm}
  \begin{columns}[T,onlytextwidth]
    \column{.47\textwidth}
      \begin{coursecallout}{coursedata!88!ink}{Valeurs de remplacement}
        \[
          \widetilde P=\MicroMedianPotions,
          \qquad \widetilde V=\MicroMedianVol.
        \]
        Ces médianes ne sont lues que sur les notes \emph{observées} du groupe
        d'estimation, jamais sur le groupe d'évaluation.
      \end{coursecallout}
    \column{.47\textwidth}
      \begin{coursecallout}{courseevaluation!88!ink}{Deux substitutions}
        \[
          \begin{aligned}
            P_{\text{Ron}}&\leftarrow\MicroMedianPotions
              &\Longrightarrow x_{\mathrm{pot}}&=0,\\
            V_{\text{George}}&\leftarrow\MicroMedianVol
              &\Longrightarrow x_{\mathrm{vol}}&=0{,}35.
          \end{aligned}
        \]
        En \PotionsName{} la médiane égale la moyenne ; en \VolName{} elle vaut
        \(\MicroMedianVol\) contre une moyenne de \(\MicroMeanVol\).
      \end{coursecallout}
  \end{columns}
  \vfill
  \centering
  \begin{adjustbox}{max width=.96\textwidth,center}
  \begin{tikzpicture}[node distance=7mm,
    box/.style={draw=ink!25,rounded corners=2pt,minimum height=11mm,
      align=center,font=\small,inner xsep=4mm},
    arr/.style={-{Latex[length=1.8mm]},thick,ink!58}]
    \node[box,fill=warning!12] (raw) {note manquante};
    \node[box,fill=coursedata!16,right=of raw] (rule) {médiane\\de la matière};
    \node[box,fill=rulecolor!16,right=of rule] (ready) {valeur\\numérique $v$};
    \node[box,fill=courseprobability!14,right=of ready] (standard)
      {$x=(v-\mu)/s$\\note standardisée};
    \draw[arr] (raw) -- (rule);
    \draw[arr] (rule) -- (ready);
    \draw[arr] (ready) -- (standard);
  \end{tikzpicture}
  \end{adjustbox}
\end{frame}

\begin{frame}[t]{Deux matières, deux échelles de mesure}
  \vspace{2mm}
  \begin{columns}[T,onlytextwidth]
    \column{.46\textwidth}
      \centering
      {\bfseries\PotionsName}
      \begin{tikzpicture}[y=.045cm]
        \draw[-{Latex[length=1.7mm]},ink!58] (0,0) -- (0,23);
        \foreach \v in {0,5,10,15,20}
          \draw[ink!55] (-.12,\v) -- (.12,\v)
            node[right=2mm,font=\scriptsize] {\v};
        \fill[coursepotions!62] (-.35,\MicroPotionsMin) rectangle (.35,\MicroPotionsMax);
      \end{tikzpicture}
      \par notée sur 20
    \column{.46\textwidth}
      \centering
      {\bfseries\VolName}
      \begin{tikzpicture}[y=.009cm]
        \draw[-{Latex[length=1.7mm]},ink!58] (0,0) -- (0,113);
        \foreach \v in {0,25,50,75,100}
          \draw[ink!55] (-.12,\v) -- (.12,\v)
            node[right=2mm,font=\scriptsize] {\v};
        \fill[coursevol!55] (-.35,\MicroVolMin) rectangle (.35,\MicroVolMax);
      \end{tikzpicture}
      \par notée sur 100
  \end{columns}
  \vfill
  Un point représente 5\% de l'échelle en \PotionsName, contre 1\% en \VolName.
\end{frame}

\begin{frame}[t]{Centrage-réduction des deux notes}
  \vspace{2mm}
  Les paramètres de standardisation sont calculés sur le groupe d'estimation,
  puis conservés :
  \begin{columns}[T,onlytextwidth]
    \column{.47\textwidth}
      \begin{coursecallout}{coursepotions!92!ink}{\PotionsName}
        \[
          \mu_{\mathrm{pot}}=\MicroMeanPotions,
          \quad s_{\mathrm{pot}}=\MicroSdPotions,
        \]
        \[
          x_{\mathrm{pot}}=
          \frac{P-\MicroMeanPotions}{\MicroSdPotions}.
        \]
      \end{coursecallout}
    \column{.47\textwidth}
      \begin{coursecallout}{coursevol!92!ink}{\VolName}
        \[
          \mu_{\mathrm{vol}}=\MicroMeanVol,
          \quad s_{\mathrm{vol}}=\MicroSdVol,
        \]
        \[
          x_{\mathrm{vol}}=
          \frac{V-\MicroMeanVol}{\MicroSdVol}.
        \]
      \end{coursecallout}
  \end{columns}
  \vfill
  Pour Harry : \(x_{\mathrm{pot}}=\MicroHarryPotions\) et
  \(x_{\mathrm{vol}}=\MicroHarryVol\). Les deux coordonnées sont désormais sans
  unité.
  \par\smallskip
  La covariance d'estimation vaut \(\MicroTrainCovariance\), soit une
  corrélation de \(\MicroTrainCorrelation\) : les deux notes varient en sens
  opposé.
\end{frame}

\begin{frame}[t]{Effet géométrique de la standardisation}
  \centering
  \begin{columns}[T,onlytextwidth]
    \column{.49\textwidth}
      \centering\small Notes brutes
      \begin{tikzpicture}
        \begin{axis}[smallcloud]
          \addplot[only marks,mark=square*,mark size=2.2pt,coursescore]
            coordinates {\MicroTrainGryffondorCoords};
          \addplot[only marks,mark=*,mark size=2.2pt,warning]
            coordinates {\MicroTrainSerpentardCoords};
        \end{axis}
      \end{tikzpicture}
    \column{.49\textwidth}
      \centering\small Coordonnées standardisées
      \begin{tikzpicture}
        \begin{axis}[stdcloud]
          \addplot[only marks,mark=square*,mark size=2.2pt,coursescore]
            coordinates {\MicroTrainGryffondorStdCoords};
          \addplot[only marks,mark=*,mark size=2.2pt,warning]
            coordinates {\MicroTrainSerpentardStdCoords};
        \end{axis}
      \end{tikzpicture}
  \end{columns}
  \vfill
  Le centrage-réduction change les unités, pas l'identité des élèves ni leur maison.
\end{frame}

% ---------------------------------------------------------------------------
\coursepart{Construire le score linéaire}{Score linéaire}{coursescore}

\begin{frame}[t]{Définition du score linéaire}
  \vspace{5mm}
  \[
    \boxed{z_i=w_0+w_{\mathrm{pot}}x_{i,\mathrm{pot}}
                    +w_{\mathrm{vol}}x_{i,\mathrm{vol}}}
  \]
  \vspace{4mm}
  \begin{columns}[T,onlytextwidth]
    \column{.47\textwidth}
      \begin{coursecallout}{coursescore!88!ink}{Coefficients}
        Le signe indique le sens de l'association ; la valeur absolue en mesure
        l'intensité sur l'échelle standardisée.
      \end{coursecallout}
    \column{.47\textwidth}
      \begin{coursecallout}{coursegradient!88!ink}{État du modèle}
        Les coefficients sont encore inconnus. Leur estimation viendra après
        la définition de la perte et du gradient.
      \end{coursecallout}
  \end{columns}
\end{frame}

\begin{frame}[t]{Contributions des deux matières au score}
  \centering
  Prenons provisoirement \(w^{\mathrm{essai}}=\MicroTrialWeights\), donc
  \(\MicroTrialScore\).
  \vspace{3mm}
  \scriptsize
  \renewcommand{\arraystretch}{1.22}
  \begin{tabular}{lrrrrrrr}
    \toprule
    élève & $x_{\mathrm{pot}}$ & $x_{\mathrm{vol}}$ &
      $-2x_{\mathrm{pot}}$ & $+x_{\mathrm{vol}}$ & $z$ & $p$ & $\ell$\\
    \midrule
    \MicroTrialRows
    \bottomrule
  \end{tabular}
  \vfill
  \normalsize
  Le score agrège les contributions ; il ne constitue encore ni une
  probabilité ni une classe prédite.
\end{frame}

\begin{frame}[t]{Le niveau de score nul définit une droite}
  \begin{columns}[c,onlytextwidth]
    \column{.60\textwidth}
      \begin{tikzpicture}
        \begin{axis}[rawcloud,legend columns=2]
          \addplot[only marks,mark=square*,mark size=2.6pt,coursescore!88!ink]
            coordinates {\MicroTrainGryffondorCoords};
          \addlegendentry{\GryffondorName}
          \addplot[only marks,mark=*,mark size=2.6pt,warning!88!ink]
            coordinates {\MicroTrainSerpentardCoords};
          \addlegendentry{\SerpentardName}
          \microboundary{ink!78,very thick};
          \node[font=\scriptsize,text=coursescore!88!ink,anchor=south]
            at (axis cs:4,74) {$z>0$};
          \node[font=\scriptsize,text=warning!88!ink,anchor=north]
            at (axis cs:16,24) {$z<0$};
        \end{axis}
      \end{tikzpicture}
    \column{.36\textwidth}
      \begin{coursecallout}{coursescore!88!ink}{Une droite, deux demi-plans}
        L'ensemble \(z=0\) est une droite du plan des notes ; elle le partage
        en deux demi-plans où le score est de signe constant.
      \end{coursecallout}
      \par\smallskip
      {\small Après estimation des coefficients, son équation sera
      \(\MicroBoundaryRaw\).}
  \end{columns}
\end{frame}

\begin{frame}[t]{Écriture matricielle des scores}
  Pour les \MicroNTrain{} observations du groupe d'estimation :
  \[
    X=
    \begin{pmatrix}
      1 & x_{1,\mathrm{pot}} & x_{1,\mathrm{vol}}\\
      \vdots & \vdots & \vdots\\
      1 & x_{\MicroNTrain,\mathrm{pot}} & x_{\MicroNTrain,\mathrm{vol}}
    \end{pmatrix},
    \qquad
    w=\begin{pmatrix}w_0\\w_{\mathrm{pot}}\\w_{\mathrm{vol}}\end{pmatrix}.
  \]
  \[
    \boxed{z=Xw}
  \]
  \vfill
  Une même multiplication produit les \MicroNTrain{} scores ; chaque ligne de
  \(X\) correspond à un élève.
\end{frame}

% ---------------------------------------------------------------------------
\coursepart{Transformer le score en probabilité}{Sigmoïde}{courseprobability}

\begin{frame}[t]{La sigmoïde associe une probabilité à tout score réel}
  \centering
  \begin{tikzpicture}
    \begin{axis}[
      width=.78\textwidth,height=5.5cm,
      xmin=-5,xmax=5,ymin=-.05,ymax=1.05,
      axis lines=middle,grid=major,grid style={gridgray!35},
      xlabel={$z$},ylabel={$p$},
      xtick={-4,-2,0,2,4},ytick={0,.5,1},
      tick label style={font=\scriptsize},label style={font=\scriptsize}]
      \addplot[courseprobability!88!ink,very thick,domain=-5:5,samples=120]
        {1/(1+exp(-x))};
      \addplot[ink!45,dashed,domain=-5:5,samples=2] {.5};
      \addplot[only marks,mark=*,mark size=2.5pt,coursescore]
        coordinates {(0,.5)};
    \end{axis}
  \end{tikzpicture}
  \[
    \boxed{p_i=\sigma(z_i)=\frac{1}{1+e^{-z_i}}}
  \]
\end{frame}

\begin{frame}[t]{Propriétés de la fonction logistique}
  \vspace{4mm}
  \begin{columns}[T,onlytextwidth]
    \column{.31\textwidth}
      \begin{coursecallout}{courseprobability!88!ink}{Monotonie}
        \[
          z_1<z_2\Rightarrow\sigma(z_1)<\sigma(z_2).
        \]
      \end{coursecallout}
    \column{.31\textwidth}
      \begin{coursecallout}{coursescore!88!ink}{Point central}
        \[
          \sigma(0)=\frac12.
        \]
      \end{coursecallout}
    \column{.31\textwidth}
      \begin{coursecallout}{courseloss!88!ink}{Symétrie}
        \[
          \sigma(-z)=1-\sigma(z).
        \]
      \end{coursecallout}
  \end{columns}
  \vfill
  \centering
  La sigmoïde conserve l'ordre des scores et les transforme en valeurs
  strictement comprises entre \(0\) et \(1\).
\end{frame}

\begin{frame}[t]{Le logit est l'inverse de la sigmoïde}
  \vspace{5mm}
  \begin{columns}[T,onlytextwidth]
    \column{.47\textwidth}
      \begin{coursecallout}{courseprobability!88!ink}{Du score à la probabilité}
        \[
          \frac{p}{1-p}=e^z,
          \qquad p=\frac{e^z}{1+e^z}.
        \]
      \end{coursecallout}
    \column{.47\textwidth}
      \begin{coursecallout}{coursescore!88!ink}{De la probabilité au score}
        \[
          \operatorname{logit}(p)
          =\ln\!\left(\frac{p}{1-p}\right)=z.
        \]
      \end{coursecallout}
  \end{columns}
  \vfill
  \centering
  \begin{tabular}{c|ccccc}
    $z$ & $-2$ & $-1$ & $0$ & $1$ & $2$\\
    \midrule
    $p$ & $0{,}119$ & $0{,}269$ & $1/2$ & $0{,}731$ & $0{,}881$
  \end{tabular}
\end{frame}

\begin{frame}[t]{Probabilités associées aux scores d'essai}
  \centering
  \renewcommand{\arraystretch}{1.35}
  \begin{tabular}{llccl}
    \toprule
    élève & maison observée & $z$ d'essai & $p=\sigma(z)$ & maison favorisée\\
    \midrule
    \MicroTrialProbabilityRows
    \bottomrule
  \end{tabular}
  \vfill
  \begin{coursecallout}{courseprobability!88!ink}{}
    La sigmoïde étant strictement croissante, deux observations de même score
    reçoivent la même probabilité selon le modèle. Les lignes marquées
    \textcolor{warning}{$\bullet$} sont celles pour lesquelles la maison
    favorisée contredit la maison observée.
  \end{coursecallout}
\end{frame}

% ---------------------------------------------------------------------------
\coursepart{Définir la perte logistique}{Perte logistique}{courseloss}

\begin{frame}[t]{Codage de la maison observée}
  \vspace{5mm}
  \[
    y_i=
    \begin{cases}
      1 & \text{si l'élève appartient à \GryffondorName},\\
      0 & \text{si l'élève appartient à \SerpentardName}.
    \end{cases}
  \]
  \vspace{5mm}
  \begin{resultat}
    Ce codage n'était nécessaire ni pour calculer le score \(z_i\), ni pour le
    transformer en probabilité \(p_i\). Il devient nécessaire pour comparer la
    sortie du modèle à la maison observée.
  \end{resultat}
\end{frame}

\begin{frame}[t]{Perte logistique d'une observation}
  \vspace{3mm}
  \[
    \boxed{\ell(z,y)=-y\ln p-(1-y)\ln(1-p)}
  \]
  \begin{columns}[T,onlytextwidth]
    \column{.47\textwidth}
      \begin{coursecallout}{coursescore!88!ink}{Maison observée : \GryffondorName}
        \[
          y=1,\qquad \ell=-\ln p=\ln(1+e^{-z}).
        \]
        La perte décroît lorsque le score augmente.
      \end{coursecallout}
    \column{.47\textwidth}
      \begin{coursecallout}{warning!88!ink}{Maison observée : \SerpentardName}
        \[
          y=0,\qquad \ell=-\ln(1-p)=\ln(1+e^z).
        \]
        La perte décroît lorsque le score diminue.
      \end{coursecallout}
  \end{columns}
\end{frame}

\begin{frame}[t]{Comportement de la perte selon la maison observée}
  \centering
  \begin{tikzpicture}
    \begin{axis}[width=.74\textwidth,height=5.2cm,xmin=-4,xmax=4,ymin=0,ymax=4.2,
      axis lines=box,grid=major,grid style={gridgray!30},
      xlabel={$z$},ylabel={$\ell$},legend style={draw=none,font=\scriptsize}]
      \addplot[coursescore!88!ink,thick,domain=-4:4,samples=100]
        {ln(1+exp(-x))};
      \addlegendentry{\GryffondorName{} observé}
      \addplot[warning!88!ink,thick,domain=-4:4,samples=100]
        {ln(1+exp(x))};
      \addlegendentry{\SerpentardName{} observé}
    \end{axis}
  \end{tikzpicture}
  \[
    \ell(z,y)=\operatorname{softplus}(z)-yz.
  \]
\end{frame}

\begin{frame}[t]{Influence des profils en recouvrement}
  Avec les coefficients d'essai, le signe du score concorde avec la maison
  observée pour \MicroTrialAgree{} élèves et s'y oppose pour
  \MicroTrialDisagree{} :
  \[
    J=\frac1{\MicroNTrain}\sum_{i=1}^{\MicroNTrain}\ell(z_i,y_i)
      =\MicroTrialLoss.
  \]
  \vspace{3mm}
  \begin{columns}[T,onlytextwidth]
    \column{.47\textwidth}
      \begin{coursecallout}{rulecolor!88!ink}{Signe concordant}
        \MicroTrialBestName{} : \(z=\MicroTrialBestScore\), \(y=1\), perte
        \(\MicroTrialBestLoss\).
      \end{coursecallout}
    \column{.47\textwidth}
      \begin{coursecallout}{warning!88!ink}{Signe discordant}
        \MicroTrialWorstName{} : \(z=\MicroTrialWorstScore\), \(y=1\), perte
        \(\MicroTrialWorstLoss\), soit un rapport de \MicroTrialLossRatio.
      \end{coursecallout}
  \end{columns}
  \vfill
  Les observations situées dans l'agrégat de la maison opposée interdisent
  toute séparation linéaire parfaite et concentrent l'essentiel de la perte.
\end{frame}

% ---------------------------------------------------------------------------
\coursepart{Calculer le gradient}{Gradient}{coursegradient}

\begin{frame}[t]{Le gradient est une moyenne de résidus pondérés}
  \vspace{4mm}
  \[
    \frac{\partial\ell_i}{\partial z_i}=p_i-y_i.
  \]
  Comme \(z_i=w_0+w_{\mathrm{pot}}x_{i,\mathrm{pot}}
                   +w_{\mathrm{vol}}x_{i,\mathrm{vol}}\),
  \[
  \boxed{
    \nabla J(w)=\frac1n\sum_{i=1}^{n}(p_i-y_i)
      \begin{pmatrix}1\\x_{i,\mathrm{pot}}\\x_{i,\mathrm{vol}}\end{pmatrix}
    =\frac1nX^\top(p-y).}
  \]
\end{frame}

\begin{frame}[t]{Gradient au point d'initialisation}
  \[
    w^{(0)}=(0,0,0)\quad\Longrightarrow\quad z_i=0,
    \quad p_i=\frac12.
  \]
  Les \MicroNTrain{} observations donnent exactement
  \[
    \sum_i(p_i-y_i)=\MicroGradientZeroSumOne,\qquad
    \sum_i(p_i-y_i)x_{i,\mathrm{pot}}=\MicroGradientZeroSumPotions,\qquad
    \sum_i(p_i-y_i)x_{i,\mathrm{vol}}=\MicroGradientZeroSumVol.
  \]
  Donc
  \[
    \boxed{\nabla J(w^{(0)})=\MicroGradientZero.}
  \]
  La première composante est non nulle : les maisons ne sont pas également
  représentées, l'ordonnée à l'origine bouge dès le premier pas.
\end{frame}

\begin{frame}[t]{Contribution de cinq observations au gradient}
  \centering
  \renewcommand{\arraystretch}{1.22}
  \begin{tabular}{lrrrrr}
    \toprule
    élève & $y$ & $p-y$ à $w=0$ & $x_{\mathrm{pot}}$ &
      $(p-y)x_{\mathrm{pot}}$ & $(p-y)x_{\mathrm{vol}}$\\
    \midrule
    \MicroResidualRows
    \bottomrule
  \end{tabular}
  \vfill
  Les contributions sont additionnées avant la division par
  \(n=\MicroNTrain\).
\end{frame}

\begin{frame}[t]{Calcul vectoriel du gradient}
  \vspace{3mm}
  \[
    \begin{aligned}
      z&=Xw,\\
      p&=\sigma(z),\\
      r&=p-y,\\
      \nabla J(w)&=\frac1{\MicroNTrain}X^\top r.
    \end{aligned}
  \]
  \vfill
  \begin{coursecallout}{coursegradient!88!ink}{Dimensions}
    \(X\in\mathbb R^{\MicroNTrain\times3}\), \(w\in\mathbb R^3\),
    \(p,y,r\in\mathbb R^{\MicroNTrain}\) et
    \(\nabla J(w)\in\mathbb R^3\).
  \end{coursecallout}
\end{frame}

\begin{frame}[t]{Vérification du gradient par différence finie}
  \[
    \frac{\partial J}{\partial w_j}\approx
      \frac{J(w+h e_j)-J(w-h e_j)}{2h},\qquad h=10^{-5}.
  \]
  \centering
  \scriptsize
  \renewcommand{\arraystretch}{1.2}
  \begin{tabular}{crrr}
    \toprule
    composante $j$ & gradient analytique & différence finie & écart\\
    \midrule
    \MicroGradientCheckRows
    \bottomrule
  \end{tabular}
  \vfill
  Écart maximal : \(\MicroGradientMaxError\).
\end{frame}

% ---------------------------------------------------------------------------
\coursepart{Estimer les coefficients}{Mise à jour}{courseupdate}

\begin{frame}[t]{Première mise à jour des coefficients}
  \[
    w^{(t+1)}=w^{(t)}-\alpha\nabla J(w^{(t)}),\qquad \alpha=1.
  \]
  À partir de \(w^{(0)}=(0,0,0)\) :
  \[
    \begin{aligned}
      w^{(1)}&=(0,0,0)-\MicroGradientZero\\
             &=\boxed{\MicroFirstStep}.
    \end{aligned}
  \]
  \vspace{1mm}
  La perte moyenne passe de \(\MicroLossAtZero\) à \(\MicroLossAfterStep\) :
  une itération réduit le coût sans encore fixer de règle de décision.
\end{frame}

\begin{frame}[t]{Convergence de la descente de gradient}
  \centering
  \scriptsize
  \renewcommand{\arraystretch}{1.08}
  \begin{tabular}{rrrrrr}
    \toprule
    itération & $J$ & $w_0$ & $w_{\mathrm{pot}}$ & $w_{\mathrm{vol}}$ &
      $\lVert\nabla J\rVert$\\
    \midrule
    \MicroGDTableRows
    \bottomrule
  \end{tabular}
  \vfill
  Le critère \(\lVert\nabla J\rVert\leq10^{-12}\) est atteint à l'itération
  \(\MicroGDIterations\).
\end{frame}

\begin{frame}[t]{Ce que caractérise l'optimum}
  L'optimum n'admet pas de forme close : il est caractérisé par l'annulation du
  gradient. Les résidus \(p_i-y_i\) y sont de somme nulle et orthogonaux aux
  deux notes.
  \par
  \centering
  \small
  \renewcommand{\arraystretch}{1.2}
  \begin{tabular}{ccc}
    \toprule
    colonne de $X$ & condition & valeur atteinte\\
    \midrule
    \MicroFirstOrderRows
    \bottomrule
  \end{tabular}
  \normalsize
  \vfill
  \begin{coursecallout}{courseupdate!88!ink}{Unicité de l'optimum}
    Newton–Raphson converge en une dizaine d'itérations, la descente de
    gradient à pas constant en \MicroGDIterations{}. Les deux méthodes
    atteignent le même vecteur : la perte logistique est strictement convexe
    sous recouvrement des deux maisons.
  \end{coursecallout}
\end{frame}

\begin{frame}[t]{Paramètres standardisés et paramètres bruts}
  \begin{columns}[T,onlytextwidth]
    \column{.48\textwidth}
      \begin{coursecallout}{courseupdate!88!ink}{Variables standardisées}
        \begin{adjustbox}{max width=\linewidth,center}
          $(w_0,w_{\mathrm{pot}},w_{\mathrm{vol}})=\MicroWDecimal$
        \end{adjustbox}
        \par\smallskip
        {\scriptsize Le coefficient de \PotionsName{} vaut deux fois celui de
        \VolName{} : un écart-type de \PotionsName{} déplace le score deux
        fois plus qu'un écart-type de \VolName.}
      \end{coursecallout}
    \column{.48\textwidth}
      \begin{coursecallout}{coursescore!88!ink}{Notes brutes}
        \begin{adjustbox}{max width=\linewidth,center}
          $(b,\beta_{\mathrm{pot}},\beta_{\mathrm{vol}})=\MicroRawBetaDecimal$
        \end{adjustbox}
        \par\smallskip
        {\scriptsize Le rapport s'inverse sur les notes brutes : un point de
        \PotionsName{} équivaut à dix points de \VolName, les deux échelles n'étant pas
        commensurables.}
      \end{coursecallout}
  \end{columns}
  \vfill
  \[
    \boxed{\MicroRawScore}
  \]
  L'ordonnée à l'origine n'est pas nulle et les deux poids n'ont pas la même
  intensité : la frontière ne passe ni par le barycentre ni le long d'une
  bissectrice. \(J(w^\star)=\MicroTrainLoss\).
\end{frame}

\begin{frame}[t]{Le recouvrement rend l'optimum fini}
  \centering
  \begin{tikzpicture}
    \begin{axis}[width=.74\textwidth,height=4.8cm,xmin=0,xmax=2.6,
      ymin=.45,ymax=.72,axis lines=box,grid=major,grid style={gridgray!30},
      xlabel={$t$ dans $w=t\,w^\star$},ylabel={$J(t)$},
      tick label style={font=\scriptsize},label style={font=\scriptsize}]
      \addplot[courseupdate!88!ink,very thick,domain=0:2.6,samples=140]
        {\MicroLossCurveExpr};
      \addplot[only marks,mark=*,mark size=2.8pt,warning]
        coordinates {(1,\MicroTrainLossPlain)};
      \node[font=\scriptsize,anchor=north west,text=warning]
        at (axis cs:1,\MicroTrainLossPlain) {$t=1$};
    \end{axis}
  \end{tikzpicture}
  \begin{releve}
    \MicroWrongSideNames{} se situent du côté opposé à leur maison. En
    l'absence de ces \MicroWrongSideCount{} observations les deux agrégats
    seraient linéairement séparables : la perte décroîtrait vers zéro lorsque
    \(t\) croît, les coefficients divergeraient et l'optimum ne serait pas
    atteint.
  \end{releve}
\end{frame}

\begin{frame}[t]{Frontière du modèle ajusté}
  \begin{columns}[c,onlytextwidth]
    \column{.60\textwidth}
      \begin{tikzpicture}
        \begin{axis}[rawcloud,legend columns=3]
          \addplot[only marks,mark=square*,mark size=2.6pt,coursescore!88!ink]
            coordinates {\MicroTrainGryffondorCoords};
          \addlegendentry{\GryffondorName}
          \addplot[only marks,mark=*,mark size=2.6pt,warning!88!ink]
            coordinates {\MicroTrainSerpentardCoords};
          \addlegendentry{\SerpentardName}
          \microboundary{ink!82,very thick};
          \addlegendentry{$p=0{,}50$}
          \MicroWrongSideLabels
        \end{axis}
      \end{tikzpicture}
    \column{.36\textwidth}
      \begin{coursecallout}{warning!88!ink}{Observations mal classées}
        \MicroWrongSideNames.
      \end{coursecallout}
      \par\smallskip
      {\small \(\MicroBoundaryRaw\)\par\smallskip
      Le modèle ajusté affecte ces \MicroWrongSideCount{} observations à la
      maison opposée.}
  \end{columns}
\end{frame}

% ---------------------------------------------------------------------------
\coursepart{Définir la règle de décision}{Décision}{coursedecision}

\begin{frame}[t]{De la probabilité à la maison prédite}
  \vspace{4mm}
  \[
    \widehat y_i=\mathbf1[p_i\geq\tau].
  \]
  \begin{columns}[T,onlytextwidth]
    \column{.47\textwidth}
      \begin{coursecallout}{warning!88!ink}{$p_i<\tau$}
        Maison prédite : \SerpentardName.
      \end{coursecallout}
    \column{.47\textwidth}
      \begin{coursecallout}{coursescore!88!ink}{$p_i\geq\tau$}
        Maison prédite : \GryffondorName.
      \end{coursecallout}
  \end{columns}
  \vfill
  Le seuil agit après l'estimation : le modifier ne recalcule ni la
  standardisation, ni les coefficients.
\end{frame}

\begin{frame}[t]{Distribution de la perte sur le groupe d'estimation}
  \[
    \ell_i=-\ln p_i \text{ si \GryffondorName},
    \qquad
    \ell_i=-\ln(1-p_i) \text{ si \SerpentardName}.
  \]
  \vspace{2mm}
  \begin{columns}[T,onlytextwidth]
    \column{.31\textwidth}
      \begin{coursecallout}{rulecolor!88!ink}{Au centre de son agrégat}
        \MicroCheapName{} : \(p=\MicroCheapProbability\), perte
        \(\MicroCheapLoss\).
      \end{coursecallout}
    \column{.31\textwidth}
      \begin{coursecallout}{coursedecision!88!ink}{Au voisinage de la frontière}
        \MicroBorderName{} : \(p=\MicroBorderProbability\), perte
        \(\MicroBorderLoss\).
      \end{coursecallout}
    \column{.31\textwidth}
      \begin{coursecallout}{warning!88!ink}{Dans l'agrégat opposé}
        \MicroExpensiveName{} : \(p=\MicroExpensiveProbability\), perte
        \(\MicroExpensiveLoss\).
      \end{coursecallout}
  \end{columns}
  \vfill
  \centering
  La perte moyenne \(J(w^\star)=\MicroTrainLoss\) est concentrée sur un petit
  nombre d'observations : elle mesure le coût du recouvrement.
\end{frame}

\begin{frame}[t]{Résultats sur le groupe d'estimation}
  \centering
  \scriptsize
  \renewcommand{\arraystretch}{.96}
  \begin{tabular}{cllrrrr}
    \toprule
    ID & élève & maison obs. & score & prob. & préd. 0,50 & préd. 0,65\\
    \midrule
    \MicroTrainCalcRows
    \bottomrule
  \end{tabular}
\end{frame}

\begin{frame}[t]{Matrice de confusion sur le groupe d'estimation}
  \centering
  \begin{tikzpicture}[x=1.45cm,y=1.02cm,
    cell/.style={draw=white,line width=1pt,minimum width=2.8cm,
      minimum height=1.35cm,font=\Large\bfseries},
    lab/.style={font=\small\bfseries,text=ink!75}]
    \node[lab,align=center] at (1,3.35) {\SerpentardName\\prédit};
    \node[lab,align=center] at (3,3.35) {\GryffondorName\\prédit};
    \node[lab,align=right,anchor=east] at (-.05,2.25) {\SerpentardName\\observé};
    \node[lab,align=right,anchor=east] at (-.05,.75) {\GryffondorName\\observé};
    \node[cell,fill=rulecolor!18,text=rulecolor!82!ink] at (1,2.25)
      {VN $=\MicroTrainTN$};
    \node[cell,fill=warning!18,text=warning!88!ink] at (3,2.25)
      {FP $=\MicroTrainFP$};
    \node[cell,fill=warning!18,text=warning!88!ink] at (1,.75)
      {FN $=\MicroTrainFN$};
    \node[cell,fill=rulecolor!18,text=rulecolor!82!ink] at (3,.75)
      {VP $=\MicroTrainTP$};
  \end{tikzpicture}
  \vfill
  \[
    \text{exactitude}
      =\frac{\MicroTrainTP+\MicroTrainTN}{\MicroNTrain}=\MicroTrainAccuracy,
    \qquad
    \text{prédire toujours \GryffondorName}
      =\frac{\MicroTrainGryffondor}{\MicroNTrain}=\MicroTrainBaseline.
  \]
\end{frame}

\begin{frame}[t]{Le seuil déplace la frontière sans réestimer le modèle}
  \begin{columns}[c,onlytextwidth]
    \column{.60\textwidth}
      \begin{tikzpicture}
        \begin{axis}[rawcloud,legend columns=2]
          \addplot[only marks,mark=square*,mark size=2.2pt,coursescore,forget plot]
            coordinates {\MicroTrainGryffondorCoords};
          \addplot[only marks,mark=*,mark size=2.2pt,warning,forget plot]
            coordinates {\MicroTrainSerpentardCoords};
          \microboundary{ink!80,thick};
          \addlegendentry{$\tau=0{,}50$}
          \microboundaryhigh{coursedecision!88!ink,dashed,thick};
          \addlegendentry{$\tau=0{,}65$}
        \end{axis}
      \end{tikzpicture}
    \column{.36\textwidth}
      \begin{coursecallout}{coursedecision!88!ink}{Translation, pas rotation}
        \begin{adjustbox}{max width=\linewidth,center}
          $b+\beta_{\mathrm{pot}}P+\beta_{\mathrm{vol}}V
            \geq\operatorname{logit}(\tau)$
        \end{adjustbox}
      \end{coursecallout}
      \par\smallskip
      {\small Pente inchangée, ordonnée à l'origine relevée de
      \MicroBoundaryShift{} points de \VolName{} : relever le seuil à
      \(0{,}65\) revient à exiger un \VolName{} supérieur à niveau de
      \PotionsName{} constant.}
  \end{columns}
\end{frame}

% ---------------------------------------------------------------------------
\coursepart{Évaluer le modèle hors échantillon}{Évaluation}{courseevaluation}

\begin{frame}[t]{Les paramètres sont figés avant l'évaluation}
  \vspace{5mm}
  \centering
  \begin{adjustbox}{max width=\textwidth,center}
  \begin{tikzpicture}[node distance=5mm,
    box/.style={draw=ink!25,rounded corners=2pt,minimum height=16mm,
      align=center,font=\scriptsize,inner xsep=3mm},
    arr/.style={-{Latex[length=2mm]},thick,ink!60}]
    \node[box,fill=courseevaluation!14] (row) {une observation\\à évaluer};
    \node[box,fill=coursedata!15,right=of row] (impute)
      {$\widetilde P=\MicroMedianPotions$ ; $\widetilde V=\MicroMedianVol$\\si une note manque};
    \node[box,fill=courseprobability!15,right=of impute] (scale)
      {$\mu=(\MicroMeanPotions,\MicroMeanVol)$\\$s=(\MicroSdPotions,\MicroSdVol)$};
    \node[box,fill=courseupdate!15,right=of scale] (model)
      {$\MicroRawScore$\\coefficients figés};
    \node[box,fill=courseevaluation!20,right=of model] (out)
      {$p$, maison prédite, perte};
    \draw[arr] (row) -- (impute);
    \draw[arr] (impute) -- (scale);
    \draw[arr] (scale) -- (model);
    \draw[arr] (model) -- (out);
  \end{tikzpicture}
  \end{adjustbox}
  \vfill
  Aucune statistique et aucun coefficient ne sont réestimés avec les huit élèves.
\end{frame}

\begin{frame}[t]{Prédictions sur le groupe d'évaluation}
  \centering
  \scriptsize
  \renewcommand{\arraystretch}{1.11}
  \begin{tabular}{cllrrrr}
    \toprule
    ID & élève & maison obs. & score & prob. & préd. 0,50 & préd. 0,65\\
    \midrule
    \MicroTestCalcRows
    \bottomrule
  \end{tabular}
  \vfill
  Perte moyenne : \(J_{\mathrm{eval}}=\MicroTestLoss\).
\end{frame}

\begin{frame}[t]{Analyse locale autour du seuil de décision}
  \begin{columns}[c,onlytextwidth]
    \column{.60\textwidth}
      \begin{tikzpicture}
        \begin{axis}[rawcloud,legend columns=2]
          \addplot[only marks,mark=square,mark size=3pt,very thick,coursescore!88!ink]
            coordinates {\MicroTestGryffondorCoords};
          \addlegendentry{\GryffondorName}
          \addplot[only marks,mark=o,mark size=3pt,very thick,warning!88!ink]
            coordinates {\MicroTestSerpentardCoords};
          \addlegendentry{\SerpentardName}
          \microboundary{ink!78,very thick};
          \addlegendentry{$\tau=0{,}50$}
          \microboundaryhigh{coursedecision!88!ink,dashed,thick};
          \addlegendentry{$\tau=0{,}65$}
          \MicroThresholdBandLabels
        \end{axis}
      \end{tikzpicture}
    \column{.36\textwidth}
      \begin{coursecallout}{coursedecision!88!ink}{Entre les deux droites}
        Olivier : \(p=\MicroProbabilityOlivier\).\par
        Daphné : \(p=\MicroProbabilityDaphne\).
      \end{coursecallout}
      \par\smallskip
      {\small Seules ces deux observations changent de classe prédite lors du
      passage de \(\tau=0{,}50\) à \(\tau=0{,}65\).}
  \end{columns}
\end{frame}

\begin{frame}[t]{Effet du seuil sur la matrice de confusion}
  \begin{columns}[T,onlytextwidth]
    \column{.48\textwidth}
      \centering
      {\bfseries\color{coursedecision!88!ink} $\tau=0{,}50$}\par\medskip
      \begin{tabular}{c|cc}
        & Serp. prédit & Gryff. prédit\\\hline
        Serp. obs. & $\MicroTestTN$ & \cellcolor{warning!18}$\MicroTestFP$\\
        Gryff. obs. & \cellcolor{warning!18}$\MicroTestFN$ &
          \cellcolor{rulecolor!18}$\MicroTestTP$
      \end{tabular}
      \par\medskip
      Daphné est un faux positif, Alicia un faux négatif.
    \column{.48\textwidth}
      \centering
      {\bfseries\color{coursedecision!88!ink} $\tau=0{,}65$}\par\medskip
      \begin{tabular}{c|cc}
        & Serp. prédit & Gryff. prédit\\\hline
        Serp. obs. & \cellcolor{rulecolor!18}$\MicroTestTNHigh$ &
          $\MicroTestFPHigh$\\
        Gryff. obs. & \cellcolor{warning!18}$\MicroTestFNHigh$ &
          $\MicroTestTPHigh$
      \end{tabular}
      \par\medskip
      Plus aucun faux positif ; Olivier devient un faux négatif.
  \end{columns}
  \vfill
  L'exactitude reste \(\MicroTestCorrect/\MicroNTest\) ; la nature de l'erreur
  change. Le seuil élevé achète de la précision
  (\(\MicroTestPrecisionTauHigh\)) au prix du rappel
  (\(\MicroTestRecallTauHigh\)).
\end{frame}

\begin{frame}[t]{Comparaison des métriques selon le seuil}
  \centering
  \renewcommand{\arraystretch}{1.35}
  \begin{tabular}{rrrrrr}
    \toprule
    seuil & exactitude & précision & rappel & spécificité & F1\\
    \midrule
    \MicroTestSummaryRows
    \bottomrule
  \end{tabular}
  \vfill
  À \(\tau=0{,}50\) : précision \(=\MicroTestPrecision\),
  rappel \(=\MicroTestRecall\), spécificité \(=\MicroTestSpecificity\),
  \(F_1=\MicroTestFScore\).
\end{frame}

\begin{frame}[t]{Incertitude d'évaluation sur \MicroNTest{} observations}
  \vspace{3mm}
  \[
    \widehat{\mathrm{acc}}
      =\frac{\MicroTestCorrect}{\MicroNTest}=\MicroTestAccuracy.
  \]
  L'intervalle de Wilson à 95\% vaut
  \[
    \boxed{[\MicroTestWilsonLow\,;\,\MicroTestWilsonHigh]}.
  \]
  \begin{vigilance}[Portée de l'étude de cas]
    Les \MicroTotalStudents{} observations sont synthétiques et choisies pour
    rendre chaque transformation vérifiable à la main. Elles illustrent une
    méthode ; elles ne mesurent pas la performance attendue sur une population
    réelle.
  \end{vigilance}
\end{frame}



\courseintro
\end{document}
